% ======================================================================
% HEAT article: Zero-parameter redshift evolution of late-type galaxy
%               sizes from a_0(z) \propto H(z)
%
% Target: journal article (aastex631 twocolumn)
% ======================================================================
\documentclass[twocolumn,letterpaper]{aastex631}

\usepackage{amsmath,amssymb}
\usepackage{natbib}
\usepackage{graphicx}
\usepackage{hyperref}

\newcommand{\azero}{a_0}
\newcommand{\Msun}{M_\odot}
\newcommand{\kms}{\,\mathrm{km\,s^{-1}}}
\newcommand{\lcdm}{\Lambda\mathrm{CDM}}
\newcommand{\heat}{\mathrm{HEAT}}
\newcommand{\azerolocal}{a_{0,0}}

% Suppress AASTeX default "Draft version ... / Typeset using LaTeX ... in AASTeX631" banner
% (the line is produced by \frontmatter@title@above in aastex631.cls; emptying it is the recommended fix)
\makeatletter
\renewcommand{\frontmatter@title@above}{}
\makeatother

\shorttitle{Zero-parameter galaxy size evolution from $\azero(z)\propto H(z)$}
\shortauthors{\textsc{Bogdan Cosmin Enciu}}

\begin{document}

\title{Zero-parameter redshift evolution of late-type galaxy sizes
       from $\azero(z)\propto H(z)$}

\author[0009-0008-7329-7200]{Bogdan Cosmin Enciu}
\affiliation{Independent Researcher}
\affiliation{M.Sc.\ in Computer Science, University Politehnica of Bucharest}
\email{bogdanenciu.author@gmail.com}

% ======================================================================
\begin{abstract}
We derive and check the current consistency of a
\emph{single closed-form, parameter-free observable}
for the late-type size--redshift relation at fixed baryonic
mass,
\[
  \frac{R_{\heat}(z)}{R_{0}}
   \;=\;(1+z)^{-3/4}\,\times\,
   \left[\frac{\Omega_{m}(z)}{\Omega_{m}(0)}\right]^{1/4},
\]
in which exponent and normalisation enter as separately testable
factors.  This primary prediction uses only the redshift
\emph{scaling}
\[
  \azero(z)=\azerolocal\,\frac{H(z)}{H_{0}},
\]
not the absolute value of $\azerolocal$.  Within the
modified-inertia family $\azero\propto H^{n}$, the late-type size
exponent measured by \citet{vanderWel2014} selects
$n=1.00\pm 0.07$ under the fixed-mass, deep-MOND assumptions used
below.  We call the resulting phenomenological scaling
hypothesis Hubble-Emergent Acceleration Theory (HEAT).
(i)~The first factor is the matter-era asymptote, matching the
empirical late-type size--redshift exponent measured in
thousands of galaxies by \citet{vanderWel2014} and confirmed at
higher redshift by \citet{Allen2017}.  (ii)~The second factor
is the closed-form normalisation correction, bounded and
monotonic, rising from unity at $z=0$ to
$\Omega_{m}^{-1/4}\!\approx\!1.335$ in the matter era and fixed
entirely by the Planck expansion history.  Current data are
consistent with $\mathcal{R}(z)$ at the $\sim 20$--$30\%$
precision set by local size-calibration systematics but do not
yet resolve its distinctive rise above the matter-era power
law; homogeneous Euclid and Roman samples will reduce the
per-bin uncertainty to the percent level, at which point
$\mathcal{R}(z)$ becomes a quantitative test of the
$\azero\propto H(z)$ coupling within the $\lcdm$ expansion
history.  As an internal consistency check (not a statistical
test, given $N\!=\!4$), the same scaling is jointly consistent
with the sample-mean flat-rotation velocity
($\langle v_{\heat}/v_{\rm obs}\rangle=1.00\pm0.08$) and size
compactification ($\langle R_{\rm obs}/R_{0}\rangle=0.38$
observed vs.\ HEAT prediction $0.37$) of four ALMA disks at
$z\!\sim\!4.5$, and reduces to standard MOND at $z=0$ with
SPARC $\chi^{2}$ statistically indistinguishable from
constant-$\azero$ MOND ($p=0.46$; $171$ galaxies) when the common
local normalisation $\azero(0)\simeq cH_0/(2\pi)$ is adopted.
This Milgrom--Hubble normalisation is a secondary, widely used
absolute anchor; it is not required for the headline size test.
Two new MUSE-DARK measurements directly probe the
$\azero(z)$ coupling: \citet{Ciocan2026RAR,Ciocan2026Halo}
establish a highly significant ($\gg 5\sigma$) positive
$\azero(z)$ trend over $0.33\!<\!z\!<\!1.44$ whose shape matches
the predicted $H(z)$ scaling.  On the clean $N\!=\!7$ cohort
(SPARC and Varasteanu+2025 at low $z$; the four Ciocan bin
centres at $z\!\sim\!0.55$--$1.30$ inside the
$0.33\!<\!z\!<\!1.44$ MUSE-DARK sample) the free-$K$ HEAT
shape achieves
$\chi^{2}\!=\!5.9$ for $6$ dof; the constant-$\azero$ MOND
baseline is excluded at $\Delta\chi^{2}\!\simeq\!192$; the
likelihood-ratio test for $n\!=\!1$ vs.\ free $n$ accepts
$n\!=\!1$ at $p\!=\!0.47$; and a $3\!+\!4$ predictive
cross-validation is internally consistent in both directions
($\chi^{2}_{\rm fwd}\!=\!7.89$ vs.\ $\chi^{2}_{\rm crit}\!=\!9.49$;
$\chi^{2}_{\rm back}\!=\!3.97$ vs.\ $\chi^{2}_{\rm crit}\!=\!7.81$,
both under HEAT $n\!=\!1$).  The MOND-framework $z\!=\!0$
anchor agrees with the parameter-free HEAT value to $\sim 1\%$.
The free-$K$ residual prefactor $\beta\!=\!1.29\!\pm\!0.04$ at
Planck cosmology is jointly degenerate with $H_{0}$ and stellar
mass-to-light: it reduces to $\beta\!=\!1.06\!\pm\!0.03$ under a
uniform $+0.11$\,dex M/L revision (within the Paper~III
per-galaxy M/L uncertainty floor), although that same revision
fails the forward $3\!\to\!4$ predictive cross-validation at
$95\%$; the raw cohort with the bare HEAT shape is the more
self-consistent slice of the $(\beta,n)$ calibration envelope.  In contrast, the
lensed bTFR zero-evolution result at $z\!\approx\!1$ of
\citet{Jeanneau2026} sits $\sim 4\sigma$ below the predicted
absolute BTFR shift and is treated as the principal extant
challenge.  In the RAR comparison, the free-$K$
curve is the main shape test; the $cH_0/(2\pi)$ anchor is retained
as a secondary absolute-normalisation reference.  We present HEAT
as a falsifiable hypothesis rather than a confirmed model:
pre-registered
thresholds on the scaling exponent, the size asymptote
$\mathcal{R}_{\infty}$, and the absolute BTFR/RAR evolution
specify what upcoming surveys are required to decide the
hypothesis.
\end{abstract}

\keywords{gravitation --- galaxies: high-redshift ---
  galaxies: formation --- galaxies: structure ---
  galaxies: kinematics and dynamics --- cosmology: theory}

% ======================================================================
\section{Introduction}\label{sec:intro}
% ======================================================================

Modified Newtonian Dynamics
\citep[MOND;][]{Milgrom1983,FamaeyMcGaugh2012} describes galactic
dynamics in the low-acceleration regime with a single universal
scale $\azero\!\approx\!1.2\times10^{-10}\,\mathrm{m\,s^{-2}}$ and
predicts the radial acceleration relation
\citep[RAR;][]{McGaugh2016,Lelli2016} without dark matter.
Because $\azero$ is constant, however, MOND is silent about
cosmological evolution.  At the same time, the James Webb Space
Telescope has revealed compact, rotationally supported galaxies
at $z\!\gtrsim\!4$ \citep{Labbe2023,Carniani2024,Bunker2023}
whose structural properties motivate reopening the question of
whether $\azero$ itself carries a redshift dependence.

\citet{Milgrom1999} identified the local MOND acceleration scale
with the de~Sitter horizon temperature, $\azero=cH_0/(2\pi)$.
Motivated by this coincidence and by the broader use of
Hubble-scale acceleration links in MOND phenomenology, we adopt
the testable modified-inertia scaling
\begin{equation}\label{eq:a0scale}
  \azero(z)=\azerolocal\,\frac{H(z)}{H_{0}},
\end{equation}
where $\azerolocal\equiv\azero(0)$ is a local anchor.  We call this
phenomenological model \emph{Hubble-Emergent Acceleration Theory}
(HEAT).  This article does not claim priority or a first-principles
derivation for the Hubble-scale acceleration idea itself; its
contribution is to propagate the redshift scaling into
parameter-free predictions for disk sizes, velocities, and
BTFR/RAR evolution, and to compare those predictions with current
data.  Only the secondary absolute tests require choosing the
standard Milgrom--Hubble anchor
$\azerolocal=cH_{0}/(2\pi)$.  The background cosmology is held
strictly degenerate with
$\lcdm$
by construction, so the theory is not a competing matter sector;
all deviations from standard physics are confined to the
low-acceleration regime, and the Newtonian regime is unmodified.

The idea that $\azero$ may carry a cosmological time dependence
has a long pedigree.  \citet{Milgrom1999} noted the numerical
coincidence $\azero\!\approx\!cH_{0}$ and discussed both
constant and evolving interpretations.
\citet{Hossenfelder2017} embedded a related coupling in the
covariant emergent gravity (CEG) framework, predicting a
substantially weaker redshift dependence than
Eq.~\ref{eq:a0scale}; \citet{HossenfelderMistele2018} showed
that the resulting $\azero(z)$ evolution in the RAR is the
cleanest discriminator among MOND-family frameworks.  HEAT is
distinguished from these predecessors not by the underlying idea
but by the specific observable package: a closed-form,
parameter-free size--redshift relation (Eq.~\ref{eq:Rfull}
below) with a bounded normalisation factor
(Eq.~\ref{eq:Rnorm2}), paired velocity predictions, and
pre-registered falsification thresholds tested against current
survey data.

While this article was being prepared, two new lines of
intermediate-redshift evidence relevant to the
$\azero(z)\propto H(z)$ hypothesis appeared.
\citet{Varasteanu2025} measure $\azero$ at $z\!<\!0.08$ from a
MIGHTEE-HI rotation-curve sample; \citet{Ciocan2026RAR}
(MUSE-DARK Paper III) infer $\azero(z)$ at $0.33\!<\!z\!<\!1.44$
from MUSE intermediate-redshift radial-acceleration relation
fits in three independent frameworks (DC14 uniform, best-fit
per-galaxy, and direct MOND), reporting a highly significant ($\gg 5\sigma$; Paper~III uncertainties are 95\,\% confidence intervals)
positive trend in all three; the matching halo-density rise
$\rho_{s}\propto(1\!+\!z)^{0.54\pm 0.31}$ inferred from the same
sample by \citet{Ciocan2026Halo} (Paper I) is the
DM-frame counterpart of HEAT's predicted $\azero(z)$ growth.
\citet{Jeanneau2026} (MUSE-DARK Paper II), by contrast, report
no detectable evolution of the lensed bTFR zero-point at
$z\!\sim\!1$; we treat this zero-evolution result as the most
direct extant challenge to HEAT's BTFR prediction
(\S\ref{sec:btfr}) and as a falsification target alongside the
headline size test.

This article focuses on a single falsifiable prediction chain.
In the deep-MOND equilibrium limit
$R\propto(GM_b/\azero)^{1/2}$ \citep{Milgrom1983,FamaeyMcGaugh2012},
so at fixed baryonic mass
\begin{equation}\label{eq:RofH}
  \frac{R(z)}{R_{0}}\;=\;\left[\frac{H(z)}{H_{0}}\right]^{-1/2}.
\end{equation}
We emphasise a structural property of Eq.~\ref{eq:RofH} that is
central to what follows: the ratio $R(z)/R_{0}$ depends only on
the \emph{scaling} $\azero\propto H$, not on the value of the
local anchor.  For any ansatz of the form
$\azero(z)=K\,H(z)$, the constant $K$ cancels identically when
the radius is normalised to its $z=0$ value, so the headline
size-evolution prediction is invariant under the choice of
absolute normalisation.  In particular, the common
Milgrom--Hubble anchor $\azerolocal=cH_{0}/(2\pi)$ fixes the
absolute value of $\azero$ but does \emph{not} propagate to
Eq.~\ref{eq:RofH} or its factorised form (Eq.~\ref{eq:Rfull}
below).  Absolute predictions such as flat-rotation velocities,
BTFR zero-point shifts, and the absolute level of direct RAR
measurements are therefore presented as a second layer on top of
the scaling-only size test, while the RAR \emph{shape} can also
be tested with a floated $K$.

Two consequences follow algebraically, with no free parameter at
any step.  The non-trivial assumptions entering Eq.~\ref{eq:RofH}
are three and explicit: (i) deep-MOND equilibrium scaling
$R\propto(GM_b/\azero)^{1/2}$, (ii) evaluation at fixed
baryonic mass across cosmic time, and (iii) a fixed local
size anchor $R_{0}$.  Throughout this article, condition~(ii) is
interpreted as a statement about \emph{mass-selected populations}
at each redshift, not about the evolutionary track of an
individual galaxy: Eq.~\ref{eq:RofH} predicts the median size of
disks of a given baryonic mass observed at redshift $z$, and the
empirical comparisons in \S\ref{sec:tests} are correspondingly
binned in stellar (or baryonic) mass at each $z$.  Departures
from any of these---evolving
mass profiles, active mergers, non-equilibrium kinematics, or
an incorrect $R_{0}$---contribute as \emph{scatter} around the
predicted curve rather than as a coherent shift in its functional
form, because non-equilibrium perturbations (mergers, tidal
interactions, gas inflows) act stochastically at any given redshift
bin and do not generically correlate with the background expansion
rate $H(z)$ that sets the predicted curve.  The
$-3/4$ exponent and the $\mathcal{R}(z)$ normalisation are
therefore protected against the generic galaxy-formation
systematics that dominate individual-galaxy residuals, provided
the sample is large enough that the stochastic perturbations
average down---a condition well satisfied by the
$\sim 10^{4}$-galaxy \citet{vanderWel2014} sample but not by the
$N\!=\!4$ ALMA cross-check.

\emph{First}, in the matter era
$H(z)\propto(1+z)^{3/2}$, so Eq.~\ref{eq:RofH} asymptotes to
$R(z)/R_{0}\propto(1+z)^{-3/4}$---the exponent measured by
\citet{vanderWel2014} in thousands of late-type galaxies and
confirmed at higher redshift by \citet{Allen2017}.
\emph{Second}, away from the matter asymptote the full
$[H(z)/H_{0}]^{-1/2}$ curve differs from the pure $(1+z)^{-3/4}$
power law by a closed-form normalisation factor
$\mathcal{R}(z)$, derived explicitly below in
Eq.~\ref{eq:Rnorm2}.  The two consequences combine into the
single composite relation
$R_{\heat}(z)/R_{0}=(1+z)^{-3/4}\times[\Omega_{m}(z)/\Omega_{m}(0)]^{1/4}$
(Eq.~\ref{eq:Rfull} below), displaying exponent and normalisation
as separate factors.  The exponent probes the matter-era
logarithmic slope, the normalisation probes the
$\Lambda$-to-matter transition shape and asymptote; both are
fixed entirely by Planck cosmology and the $\azero\propto H(z)$
coupling.

The rest of this article formalises the derivation
(\S\ref{sec:ansatz}--\S\ref{sec:predictions}), compares the
predictions to existing data (\S\ref{sec:tests}), and
pre-registers falsification thresholds (\S\ref{sec:falsify}).
We defer additional phenomenology---a baryon-budget relief
channel for JWST $z\gtrsim 10$ candidates, collapse-time
enhancements, and the full SPARC $\chi^2$ analysis---to a longer
treatment and do not advance it as primary evidence here.

% ======================================================================
\section{Scaling ansatz and optional absolute anchor}\label{sec:ansatz}
% ======================================================================

We pose the question empirically rather than thermodynamically.
Within the modified-inertia family
$\azero(z)\propto H^{n}(z)$, the deep-MOND closure
$R\propto(GM_b/\azero)^{1/2}$ \citep{Milgrom1983,FamaeyMcGaugh2012}
combined with the matter-era scaling
$H(z)\propto(1+z)^{3/2}$ forces a late-type size exponent
$d\log R/d\log(1+z)=-3n/4$.  The measured value
$-0.75\!\pm\!0.05$ from the deep \citet{vanderWel2014}
late-type sample therefore selects
$n=1.00\!\pm\!0.07$ (\S\ref{sec:exp}); within this family the
redshift evolution collapses to the scaling of Eq.~\ref{eq:a0scale}.
The role of this section is to separate the scaling-only
predictions from the optional absolute anchor used later for
velocity, BTFR, RAR, and SPARC comparisons.

\paragraph{Status of free parameters.}
Two distinct claims must be kept separate.
\emph{Redshift evolution:} once the
$n=1$ scaling $\azero(z)/\azerolocal=H(z)/H_0$ is adopted, the
size and velocity scalings of Eqs.~\ref{eq:Rfull}--\ref{eq:Vfull}
introduce no additional fitted redshift-evolution parameter.
The empirical motivation for this choice is the matter-era size
exponent (\S\ref{sec:exp}).
\emph{Absolute normalisation at $z\!=\!0$:} a separate local
anchor is required only for absolute accelerations and velocities.
For those secondary tests we use the standard Milgrom--Hubble
value $\azerolocal=cH_0/(2\pi)$, which matches the empirical MOND
scale to within $13\%$.  The headline observable $R(z)/R_{0}$
depends on the first claim only (K-invariance, \S\ref{sec:exp}).

\paragraph{Motivation for the absolute anchor.}
An observer comoving with the FRW background sees an apparent
horizon at radius $R_A(z)=c/H(z)$ with surface gravity
$\kappa_A(z)=cH(z)$ \citep{BakRey2000,CaiKim2005} and a
Hawking--Unruh temperature
$T_A=\hbar\kappa_A/(2\pi c\,k_B)=\hbar H/(2\pi k_B)$
\citep{Hawking1975,Unruh1976,GibbonsHawking1977}.  At $z=0$
\citet{Milgrom1999} observed that the corresponding acceleration
scale, $cH_0/(2\pi)\!\approx\!1.04\times10^{-10}
\,\mathrm{m\,s^{-2}}$, lies within $13\%$ of the empirical MOND
scale inferred from SPARC
($1.2\times10^{-10}\,\mathrm{m\,s^{-2}}$;
\citealt{Lelli2016,Li2018}).  When an absolute normalisation is
needed, we therefore use
\begin{equation}\label{eq:a0}
  \azero(z)\;=\;\azerolocal\,\frac{H(z)}{H_0},
  \qquad
  \azerolocal\equiv\frac{cH_0}{2\pi},
\end{equation}
which is algebraically equivalent to $\azero(z)=cH(z)/(2\pi)$.
The $2\pi$ factor is the usual Hawking--Unruh periodicity, but
we treat Eq.~\ref{eq:a0} as a phenomenological anchor rather than
as a microscopic derivation.  The point of the present paper is
not priority for the Hubble-scale acceleration link itself, but
the closed observable package obtained by applying the scaling to
disk sizes, flat velocities, and BTFR/RAR evolution.  As noted in \S\ref{sec:intro}, CEG
\citep{Hossenfelder2017,HossenfelderMistele2018} predicts a
substantially weaker redshift dependence; the two hypotheses
can be distinguished observationally (\S\ref{sec:falsify}).  The
empirical contribution of the present article is to propagate
Eq.~\ref{eq:a0scale} into zero-parameter scaling predictions and,
where an absolute value is required, to test the conventional
anchor of Eq.~\ref{eq:a0} against current redshift data.

% ======================================================================
\section{Scaling-only predictions for sizes and velocities}\label{sec:predictions}
% ======================================================================

\subsection{Derivation of the exponent}\label{sec:exp}

Once the scaling of Eq.~\ref{eq:a0scale} is adopted, the first
observable consequence is the size evolution of a mass-selected
disk population.  In the deep-MOND regime
$R\propto(GM_b/\azero)^{1/2}$ \citep{Milgrom1983,FamaeyMcGaugh2012}.
At fixed baryonic mass this yields
Eq.~\ref{eq:RofH}.  The matter-dominated cosmological scaling
$H(z)\propto(1+z)^{3/2}$ then forces
\begin{equation}\label{eq:Rexp}
  R(z)/R_{0}\;\longrightarrow\;(1+z)^{-3/4}
  \qquad(z\gg z_{\rm eq,\Lambda}),
\end{equation}
i.e.\ the exact exponent of the late-type size--redshift
relation measured by \citet{vanderWel2014} in $3$D-HST$+$CANDELS
out to $z\!\sim\!3$ and confirmed at $z\!\sim\!3$--$5$ by
\citet{Allen2017}.  The rational exponent $-3/4$ is an
observable consequence of the scaling-only assumption, the
deep-MOND kernel, and matter-era cosmology, with no fitted
quantity in the size prediction.  Within
the modified-inertia family $\azero\propto H^{n}$, the matter-era
size exponent is $-3n/4$, so the \citet{vanderWel2014} measurement
of $-0.75\pm 0.05$ corresponds to $n=1.00\pm 0.07$ under the
fixed-mass, deep-MOND interpretation.  In that restricted family,
the constant-$\azero$ limit ($n=0$, standard MOND; predicted
exponent~$0$), the fractional coupling $n=1/2$ (exponent $-3/8$),
and the quadratic coupling $n=2$ (exponent $-3/2$) do not match
the observed slope without adding additional galaxy-formation
evolution.  This is why the paper adopts $n=1$ as the benchmark
scaling rather than fitting a redshift exponent later.
To be explicit about the logical status of this choice: the
matter-era exponent $-3/4$ is the \emph{empirical motivation}
for selecting $n\!=\!1$, not the test of the resulting model.
The independent, falsifiable predictions are the
normalisation curve $\mathcal{R}(z)$ and its asymptote
$\Omega_{m}^{-1/4}$ (\S\ref{sec:norm}), the velocity scaling
(\S\ref{sec:family}), the BTFR zero-point evolution
(\S\ref{sec:btfr}), and the direct $\azero(z)$ values
(\S\ref{sec:a0z})---none of which are inputs to the
$n$-selection step.

For completeness we note that a baseline $\lcdm$ disk-formation
calculation with
$R_{\rm disk}\propto\lambda\,R_{\rm vir}\propto H^{-2/3}$
\citep{MoMaoWhite1998} returns the matter-era exponent $-1$,
while abundance-matched fits invoking halo-concentration
evolution, adiabatic contraction, and feedback-driven angular
momentum loss return best-fit exponents in the range
$-0.6$ to $-1.0$ depending on ingredient choices
\citep{Somerville2018}.  HEAT's $-3/4$ is (a)~a rational
exponent following from the adopted acceleration scaling and the
deep-MOND kernel rather than an ensemble fit, and (b)~within
$1\sigma$ of the measured value
\citep{vanderWel2014}.

\subsection{Derivation of the
  normalisation}\label{sec:norm}

The matter-era power law $(1+z)^{-3/4}$ is an asymptote, not the
full HEAT prediction.  The ratio of the full
$[H(z)/H_{0}]^{-1/2}$ curve to that asymptote defines
$\mathcal{R}(z)\equiv (1+z)^{3/4}\,[H(z)/H_{0}]^{-1/2}$ and
admits a closed form.  Writing
$H(z)^{2}=H_{0}^{2}[\Omega_{m}(1+z)^{3}+\Omega_{\Lambda}]$ and
using the identity
$\Omega_{m}(z)\equiv\Omega_{m}(1+z)^{3}/E(z)^{2}$,
\begin{equation}\label{eq:Rnorm2}
  \mathcal{R}(z)\;=\;
    \left[\frac{\Omega_{m}(z)}{\Omega_{m}(0)}\right]^{1/4}.
\end{equation}
$\mathcal{R}(z)$ is bounded, monotonic, and fully parameter-free:
$\mathcal{R}(0)=1$ by the local anchor, and
$\mathcal{R}(z)\to\Omega_{m}^{-1/4}\!\approx\!1.335$ in the
matter era.  The transition shape between the two limits is the
$\Lambda$-to-matter interpolation around
$z_{\rm eq,\Lambda}\!\approx\!0.3$, with no galaxy-formation
freedom anywhere in the curve.  Numerically, with Planck-fixed
$\Omega_m=0.315$,
\begin{equation*}
  \mathcal{R}(z)\;=\;
    \{1.00,\,1.18,\,1.25,\,1.29,\,1.31,\,1.32,\,1.33\}
\end{equation*}
at $z=\{0,\,0.5,\,0.9,\,1.5,\,2.0,\,3.0,\,4.0\}$.

Collecting the two observable consequences, Eqs.~\ref{eq:RofH} and
\ref{eq:Rnorm2} together imply a single composite prediction
for the late-type size--redshift relation,
\begin{equation}\label{eq:Rfull}
  \boxed{\;\;\frac{R_{\heat}(z)}{R_{0}}
   \;=\;(1+z)^{-3/4}\;\times\;
   \left[\frac{\Omega_{m}(z)}{\Omega_{m}(0)}\right]^{1/4}\;\;}
\end{equation}
in which the first factor is the matter-era asymptote
(\S\ref{sec:exp}) and the second is the closed-form
$\Lambda$-to-matter normalisation correction (\S\ref{sec:norm}),
asymptoting to $\Omega_{m}^{-1/4}\!\approx\!1.335$.  Equation
\ref{eq:Rfull} is algebraically equivalent to Eq.~\ref{eq:RofH}
but displays the two empirically accessible faces of the
prediction---exponent and normalisation---as separate factors,
each following from the scaling-only acceleration relation, the
deep-MOND kernel, and Planck cosmology with no fitted quantity.
Together with the paired velocity scaling in \S\ref{sec:family},
Eq.~\ref{eq:Rfull} is the distinctive observational package that
upcoming Euclid, Roman, and JWST mass-complete catalogues will
test.

Eq.~\ref{eq:Rnorm2} is the \emph{core novel prediction of this
article}.  Importantly it is not a fit: the $z=0$ anchor is definitional
(both curves equal $R_{0}$ at $z=0$ by construction), and the
matter-era asymptote $\Omega_{m}^{-1/4}$ is fixed by Planck.
Only the \emph{shape} of $\mathcal{R}(z)$ at $z>0$ carries
non-trivial content, but that shape is also fully determined.
Any coupling $\azero\propto H^{n}$ with $n\neq 1$, or a matter-only
coupling $\azero\propto H_{m}(z)$, predicts a different
asymptote; the observed saturation near $1.34$ at $z\gtrsim 2$
\citep[Table~\ref{tab:interz} below and][]{vanderWel2014} is
consistent with $n=1$ and disfavours $n\neq 1$ at the
qualitative level of current data.

\paragraph{Distinguishing feature relative to galaxy-formation
physics.}
The strongest discriminant proposed in this article is the
specific two-part shape of Eq.~\ref{eq:Rnorm2}: the asymptote
$\mathcal{R}_{\infty}\!\approx\!1.34$ and the location of the
$\Lambda$-sourced break at
$z_{\rm eq,\Lambda}\!\approx\!0.3$.  These are simultaneous
signatures of the cosmological constant entering the size
relation through $\Omega_{m}(z)$.  The closed-form prediction is
therefore sharper than the matter-era exponent alone, because
standard $\lcdm$ disk-formation modelling does not generically
single out this same bounded transition curve without additional
assumptions about the redshift evolution of baryonic ingredients.
The exponent $-3/4$ alone does not
establish this: abundance-matched $\lcdm$ disk-formation models
with tuned feedback and angular-momentum loss can reach
matter-era exponents anywhere in the range $-0.6$ to $-1.0$
\citep{Somerville2018,MoMaoWhite1998}, and the measured $-0.75$
lies inside that band.  The \emph{normalisation curve} is
different in kind.  Eq.~\ref{eq:Rnorm2} is a function of
$\Omega_{m}(z)$ alone, with asymptote
$\Omega_{m}^{-1/4}\!\approx\!1.34$ and a specific
$\Lambda$-sourced transition shape centred on
$z_{\rm eq,\Lambda}\!\approx\!0.3$.  Baryonic galaxy-formation
physics---feedback efficiencies, merger rates, spin-parameter
evolution, adiabatic contraction---is not directly tied to
$\Omega_{\Lambda}$: these processes act on halo time-scales
$\propto H(z)^{-1}$ that already absorb the expansion history,
but standard implementations do not single out the
$[\Omega_{m}(z)/\Omega_{m}(0)]^{1/4}$ shape as a generic outcome.
Reproducing Eq.~\ref{eq:Rnorm2} within $\lcdm$ is therefore
\emph{not obviously a generic outcome of disk-formation
modelling}: it would require a combination of baryonic
processes to approximate the same $z$-profile, rather than
following from a single background quantity.  HEAT, by contrast,
delivers the full $\mathcal{R}(z)$ with no adjustable redshift
parameter as an algebraic
consequence of Eq.~\ref{eq:a0scale}.  Measuring $\mathcal{R}(z)$ to
percent-level accuracy over $0\!\lesssim\!z\!\lesssim\!3$
therefore probes a structure that is non-generic in $\lcdm$
galaxy-formation models without invoking a specific,
$\Omega_\Lambda$-dependent combination of baryonic ingredients.

\paragraph{Robustness to named astrophysical systematics.}
We list the three baryonic effects most likely to be invoked
against Eq.~\ref{eq:Rnorm2} and indicate why each does not
naturally reproduce its specific shape.  (a)~\emph{Evolving gas fractions
and stellar/baryonic mass selection}: gas-rich high-$z$ disks
have larger $f_{\rm gas}$ and therefore systematically lower
$M_\star/M_b$, but the corresponding bias in size at fixed
$M_\star$ is smooth, monotone, and tracks the halo time-scale
$\propto H(z)^{-1}$; it does not single out the
$\Lambda$-sourced break near $z_{\rm eq,\Lambda}\!\approx\!0.3$
that is the distinctive feature of $\mathcal{R}(z)$.
(b)~\emph{Merger-induced size growth}: dry mergers inflate sizes
but predominantly broaden the dispersion at fixed mass rather
than coherently shift the binned mean by the
$\Omega_{m}^{-1/4}\!\approx\!1.34$ factor that
Eq.~\ref{eq:Rnorm2} forces, and recent merger-rate measurements
\citep{Conselice2014,Duncan2019} show comparable rates across
$z\!\sim\!1$--$3$ rather than the
$\Lambda$-modulated profile required.
(c)~\emph{Local size-anchor calibration drift}: the
\citet{Shen2003} versus \citet{vanderWel2014} local anchors
differ by $\sim 20\%$, comparable to the $\sim 25$--$33\%$
elevation but with no redshift dependence, so it cannot mimic
the rise from $\mathcal{R}(0)=1$ to
$\mathcal{R}_{\infty}\!\approx\!1.34$.  Reproducing the joint
shape (location of break, asymptote, monotonicity) within
$\lcdm$ would therefore require these effects, or other
baryonic ingredients, to combine into a coordinated
$\Omega_\Lambda$-dependent profile rather than each contributing
as an independent broadening.

\paragraph{Quantitative mass-selection mimic.}  We make
systematic~(a) explicit.  Suppose a static-MOND universe (constant
$\azero$) in which the underlying baryonic mass at fixed
\emph{observed} stellar mass evolves as a power law,
$M_{b}(z)=M_{b}(0)\,(1+z)^{\alpha}$, capturing evolving gas
fractions and stellar/baryonic mass selection in a one-parameter
caricature.  The deep-MOND equilibrium
$R\propto(GM_{b}/\azero)^{1/2}$ then gives, at fixed observed
$M_{\star}$,
\begin{equation}\label{eq:mimic}
  \mathcal{R}_{\rm mimic}(z;\alpha)\;\equiv\;
   \frac{R(z)/R_{0}}{(1+z)^{-3/4}}
   \;=\;(1+z)^{\alpha/2}.
\end{equation}
Eq.~\ref{eq:mimic} is monotone, unbounded, and
$\Lambda$-independent.  HEAT's $\mathcal{R}(z)$ is bounded,
saturates at $\Omega_{m}^{-1/4}\!\approx\!1.335$, and has a
$\Lambda$-sourced break at
$z_{\rm eq,\Lambda}\!\approx\!0.3$.  No single $\alpha$
reproduces both shape features.  The value
$\alpha\!\approx\!0.36$ tuned to match the asymptote at
$z\!=\!4$ \emph{by construction} still under-predicts the rise
at $z\!=\!0.5$ by $\sim 10\%$ (mimic $1.08$ vs HEAT $1.18$) and
overshoots the saturation at $z\!\gtrsim\!4$ as $z\to\infty$.
Figure~\ref{fig:mass_sel} visualises the family
$\alpha\!\in\![0.0,0.5]$ together with the asymptote-tuned
mimic and the HEAT closed form: the mimics share the
$z\!=\!0$ anchor with $\mathcal{R}(z)$ but diverge from it
both at the break ($z\!\sim\!0.3$, where they rise too slowly)
and at high redshift ($z\!\gtrsim\!4$, where they fail to
saturate).  This shows that the one-parameter
power-law mass-evolution caricature cannot by itself reproduce
$\mathcal{R}(z)$ at the percent level and quantifies
systematic~(a): an $\Omega_\Lambda$-independent power-law
$M_{b}(z)$ does not reproduce the
$\Omega_{m}^{-1/4}$ asymptote without simultaneously
violating the binned-mean rise at the break.

\begin{figure}[t]
  \centering
  \includegraphics[width=\columnwidth]{figures/fig8_mass_selection.pdf}
  \caption{Diagnostic comparison between $\mathcal{R}(z)$ and a
    one-parameter mass-selection mimic.  Solid blue: HEAT closed
    form $\mathcal{R}(z)=[\Omega_{m}(z)/\Omega_{m}(0)]^{1/4}$,
    saturating at $\Omega_{m}^{-1/4}\!\approx\!1.335$ (dotted
    horizontal line) with a break centred on
    $z_{\rm eq,\Lambda}\!\approx\!0.3$ (dotted vertical line).
    Dashed coloured curves: the family of static-MOND
    mass-evolution mimics
    $\mathcal{R}_{\rm mimic}(z;\alpha)\!=\!(1+z)^{\alpha/2}$
    arising from $M_{b}(z)\!=\!M_{b}(0)(1+z)^{\alpha}$ for
    $\alpha\!\in\!\{0.0,0.1,0.2,0.3,0.4,0.5\}$.  Red dot-dashed:
    the asymptote-tuned value $\alpha\!\approx\!0.36$ that
    matches HEAT at $z\!=\!4$ by construction; it still
    under-predicts the rise at $z\!=\!0.5$ by
    $\sim 10\%$ and fails to saturate at higher redshift.  Within
    this deliberately simple one-parameter power-law family, no
    single $\alpha$ reproduces both the asymptote and the
    location/sharpness of the $\Lambda$-sourced break.
    \label{fig:mass_sel}}
\end{figure}

\subsection{Scaling family}\label{sec:family}

The same factorisation applied to any deep-MOND observable
$\mathcal{O}\propto\azero^{p}$ yields the redshift dependence
$\mathcal{O}(z)/\mathcal{O}(0)=[H(z)/H_{0}]^{p}$ with matter-era
asymptote $(1+z)^{3p/2}$ and normalisation
$[\Omega_{m}(z)/\Omega_{m}(0)]^{-p/2}$.  Each exponent below is
rational and algebraically forced.  The size and velocity rows
carry exponents of \emph{opposite} sign, which is what makes the
joint velocity--size cross-check of \S\ref{sec:alma}
structurally informative even at small $N$.

\smallskip\noindent\textbf{One scaling, four observable
consequences.}
\begin{itemize}\setlength\itemsep{2pt}
\item[$\bullet$] \textit{$R_{\rm eff}$ at fixed $M_b$:}
  $p=-1/2$, so $R\propto[H(z)/H_0]^{-1/2}$ and
  $R\propto(1+z)^{-3/4}$ in the matter era.  This is the
  primary size-evolution shape test.
\item[$\bullet$] \textit{$v_{\rm flat}$ at fixed $M_b$:}
  $p=+1/4$, so $v\propto[H(z)/H_0]^{1/4}$ and
  $v\propto(1+z)^{3/8}$ in the matter era.  This is the
  velocity consistency check.
\item[$\bullet$] \textit{BTFR zero-point:} $p=-1$, so the
  locally anchored BTFR zero-point evolves as
  $[H(z)/H_0]^{-1}$.  This tests the same coupling with an
  independent observable.
\item[$\bullet$] \textit{RAR knee:} $p=+1$, so
  $g_{\rm knee}=\azero\propto H(z)$.  This is the direct
  absolute-acceleration test.
\end{itemize}

These four channels are not equivalent under a global rescaling
of the local anchor, $\azerolocal\!\to\!K\,\azerolocal$.  Such a
rescaling changes every absolute MOND acceleration, but it cancels
from the redshift ratios $R(z)/R_{0}$, $v(z)/v_{0}$, and the
locally anchored BTFR zero-point evolution.  The size, velocity,
and BTFR-evolution rows therefore probe the \emph{shape}
$\azero(z)/\azero(0)=H(z)/H_{0}$.  The direct RAR
measurement of \S\ref{sec:a0z} tests the same $H(z)$ shape with
a floated $K$ and, separately, constrains the absolute
normalisation when the $cH_0/(2\pi)$ anchor is imposed.  The size row is tested
against \citet{vanderWel2014}, \citet{Allen2017}, and the
calibrated \citet{Allen2025} shape check; the BTFR row carries
the \citet{Ubler2017} support and the \citet{Jeanneau2026}
challenge; the RAR row is tested directly by
\citet{Ciocan2026RAR}.

Applying the factorisation of Eq.~\ref{eq:Rfull} to the
flat-rotation velocity ($p=+1/4$) gives the composite
zero-parameter prediction
\begin{equation}\label{eq:Vfull}
  \boxed{\;\;\frac{v_{\rm flat}(z)}{v_{0}}
   \;=\;(1+z)^{+3/8}\;\times\;
   \left[\frac{\Omega_{m}(z)}{\Omega_{m}(0)}\right]^{-1/8}\;\;}
\end{equation}
with $z\!=\!0$ anchor $v_{0}\!=\!(GM_{b}\azero(0))^{1/4}$, a
velocity normalisation
$\mathcal{V}(z)\!\equiv\![\Omega_{m}(z)/\Omega_{m}(0)]^{-1/8}$,
and matter-era asymptote
$\mathcal{V}_{\infty}\!=\!\Omega_{m}^{+1/8}\!\approx\!0.866$.
The opposite signs of the exponents in $\mathcal{R}(z)$ ($+1/4$)
and $\mathcal{V}(z)$ ($-1/8$) make their ratio
$\mathcal{R}(z)/\mathcal{V}(z)=[\Omega_{m}(z)/\Omega_{m}(0)]^{3/8}$
a further parameter-free cosmological probe, asymptoting to
$\Omega_{m}^{-3/8}\!\approx\!1.54$.

We stress, however, that the observational leverage of
Eq.~\ref{eq:Vfull} is substantially weaker than that of
Eq.~\ref{eq:Rfull} on current data: the predicted normalisation
shift is only $\mathcal{V}_{\infty}\!-\!1\!\approx\!-13\%$ (versus
$\mathcal{R}_{\infty}\!-\!1\!\approx\!+34\%$ for sizes), below
the $15$--$20\%$ per-galaxy systematic floor on high-$z$ rotation
velocities from inclination, beam-smearing, and
turbulent-broadening subtraction (\S\ref{sec:alma}).  The
velocity channel therefore functions here as a \emph{consistency
requirement}---the ALMA sample mean of \S\ref{sec:alma} is
compatible with Eq.~\ref{eq:Vfull}---rather than as a precision
test of $\mathcal{V}(z)$.  A definitive test of the velocity
normalisation requires mass-binned, kinematically clean samples
with $\gtrsim 50$ disks per redshift bin, which upcoming ALMA
large programmes and JWST/NIRSpec IFU campaigns are expected to
deliver.

% ======================================================================
\section{Observational tests}\label{sec:tests}
% ======================================================================

We organise the empirical comparisons that follow into an
explicit prediction hierarchy.  The \emph{primary} observable
test is the scaling-only shape and asymptote of the late-type
size--redshift relation $R(z)/R_{0}$ (\S\ref{sec:interz}), which
is supported by the largest sample ($\sim 10^{4}$ galaxies;
\citealt{vanderWel2014}) and tests the headline prediction
Eq.~\ref{eq:Rfull} directly.  The next test is the direct
MUSE-DARK RAR measurement of $\azero(z)$: there we put the
free-$K$ curve $\azero(z)=K\,cH(z)$ on the main line because it
tests the \emph{shape} independently of the local anchor, while
the $cH_0/(2\pi)$ normalisation is shown as a secondary reference.
The paired velocity scaling in Eq.~\ref{eq:Vfull} is then checked
with the available ALMA $z\!\sim\!4.5$ sample (\S\ref{sec:alma},
$N\!=\!4$), which remains a supporting consistency check rather
than a statistical test.  BTFR zero-point evolution
(\S\ref{sec:btfr}) and the SPARC $z\!=\!0$ comparison
(\S\ref{sec:sparc}) complete the absolute-normalisation layer.
Pre-registered falsification thresholds for both layers are
collected separately in \S\ref{sec:falsify}.

\subsection{Late-type size evolution at
  $z\approx 0.9$--$8$}\label{sec:interz}

We compare the composite prediction Eq.~\ref{eq:Rfull} to the
empirical late-type size--redshift fit of \citet{vanderWel2014},
extended into the ZFOURGE regime by \citet{Allen2017}, anchored
at $R_{0}\!\approx\!4\,\mathrm{kpc}$ at $\log M_{\star}=10$
\citep{Shen2003}.  The first factor of Eq.~\ref{eq:Rfull} is
tested by the exponent; the second factor is tested by the
binned-mean normalisation offset.
Table~\ref{tab:interz} lists HEAT's $R(z)/R_{0}$ and the
corresponding empirical $(1+z)^{-0.75}$ values over
$z=0.9$--$3.5$.  The two are consistent on the exponent to
within $1\%$ across the range, while the HEAT curve lies
systematically $25$--$33\%$ above the matter-era empirical
power law---consistent within the local-anchor systematic
band with the behaviour predicted by Eq.~\ref{eq:Rnorm2},
which forces
$\mathcal{R}(z\geq 3)\to\Omega_{m}^{-1/4}\!\approx\!1.34$.
The binned-mean ratio column is consistent with the
$\mathcal{R}(z)$ profile of \S\ref{sec:norm} at the
$\sim 20$--$30\%$ precision set by local size-calibration
systematics (the \citealt{Shen2003} vs.\ \citealt{vanderWel2014}
local anchors alone differ by $\sim 20\%$).

A complementary high-$z$ check is provided by the JWST/NIRCam
rest-optical size--mass relation of \citet{Allen2025}.  At fixed
$M_{\star}=5\times10^{10}\,M_{\odot}$ they report
$\log R_{e}=(-0.807\pm 0.026)\log(1+z)+(0.947\pm 0.014)$ over
$0\!<\!z\!<\!9$.  When evaluated at the same local late-type
anchor used here,
$R_{0}=4(M_{\star}/10^{10}\,M_{\odot})^{0.22}=5.70$\,kpc, the
relation tracks $[H(z)/H_{0}]^{-1/2}$ to within
$2.7$--$8.4\%$ over $z=3$--$8$, and the HEAT-implied mean local
anchor inferred from the Allen high-$z$ points is $6.00$\,kpc,
only $5.3\%$ above the adopted one.  This is a calibrated shape
check, not an independent absolute-normalisation confirmation:
the \citet{Allen2025} extrapolated $z\!=\!0$ intercept is
$10^{0.947}\!=\!8.85$\,kpc, about $55\%$ larger than the local
anchor used in this article.  The numerical comparison at
$z\!\in\!\{3,4.5,6,8\}$ is reproduced by the
\texttt{\_allen2025\_size\_calibration} block of
\texttt{publication/jwst\_early\_galaxies.py}.

We emphasise that this is a consistency check, not a precision
test of the functional form: we do not quote a significance in
$\sigma$ units because a proper likelihood on the joint
size--mass--redshift surface is not available in the literature,
and the current data neither confirm nor exclude
$\mathcal{R}(z)$ against the matter-era baseline at the level
required.  The apparent plateau of the binned-mean ratio near
$1.3$ across $z\!\in\![2,3.5]$ sits within the local-anchor
systematic band; resolving it to the percent level is the
specific task that upcoming surveys are expected to accomplish.  Figure~\ref{fig:norm} visualises both
features: panel~(a) shows the HEAT and matter-era curves sharing
a common $z\!=\!0$ anchor and diverging by $\mathcal{R}(z)$,
while panel~(b) isolates $\mathcal{R}(z)$ and its asymptote.

\begin{deluxetable}{cccc}
\tablecaption{Intermediate-redshift size evolution: zero-parameter
  HEAT prediction vs.\ the matter-era empirical $(1+z)^{-0.75}$
  fit \citep{vanderWel2014,Allen2017}, with the predicted ratio
  $\mathcal{R}(z)$ (Eq.~\ref{eq:Rnorm2}).\label{tab:interz}}
\tablewidth{0pt}
\tablehead{
  \colhead{$z$} & \colhead{$R_{\heat}/R_{0}$}
  & \colhead{$(1+z)^{-0.75}$}
  & \colhead{$\mathcal{R}(z)=R_{\heat}/[(1+z)^{-0.75}]$}
}
\startdata
$0.9$ & $0.77$ & $0.62$ & $1.25$ \\
$1.5$ & $0.65$ & $0.50$ & $1.29$ \\
$2.0$ & $0.57$ & $0.44$ & $1.31$ \\
$3.0$ & $0.47$ & $0.35$ & $1.32$ \\
$3.5$ & $0.43$ & $0.32$ & $1.33$ \\
\enddata
\tablecomments{HEAT values use the scaling
  $\azero(z)/\azero(0)=H(z)/H_0$ with Planck-fixed cosmology; no
  absolute acceleration normalisation enters.  The rightmost
  column is the closed-form prediction of Eq.~\ref{eq:Rnorm2},
  asymptoting to $\Omega_{m}^{-1/4}\!\approx\!1.335$ at
  $z\gtrsim 3$.  Consistency with the full shape of
  $\mathcal{R}(z)$ is at the $\sim 20$--$30\%$ level of local
  size-calibration systematics.}
\end{deluxetable}

\begin{figure}[t]
  \centering
  \includegraphics[width=\columnwidth]{figures/fig6_normalization.pdf}
  \caption{Scaling-only HEAT normalisation test.
    \textit{Panel (a):} predicted size evolution
    $R_{\rm eff}(z)/R_{0}$ under the composite relation
    of Eq.~\ref{eq:Rfull}.  Solid blue: HEAT model curve,
    $R/R_{0}=(1+z)^{-3/4}[\Omega_{m}(z)/\Omega_{m}(0)]^{1/4}$,
    equivalently $[H(z)/H_{0}]^{-1/2}$.  Dashed grey: the
    matter-era first factor $(1+z)^{-3/4}$ in isolation.  Dotted
    red: reference $\lcdm$ disk-formation baseline of
    \citet{MoMaoWhite1998}, $R/R_{0}=[H(z)/H_{0}]^{-2/3}$,
    asymptoting to $(1+z)^{-1}$ in the matter era; it sits
    systematically below the HEAT curve and crosses below the
    empirical \citet{vanderWel2014} matter-era power law for
    $z\!\gtrsim\!3$, illustrating the well-known under-prediction
    of disk sizes at high redshift in the simplest
    \citet{MoMaoWhite1998} closure without
    feedback/spin/concentration tuning.  The HEAT and matter-era curves
    share the $z\!=\!0$ anchor and the high-$z$ logarithmic
    slope but differ by the closed-form second factor
    $\mathcal{R}(z)$.  The shaded band marks a
    $\pm 20\%$ envelope reflecting the local anchor systematic
    (\citealt{Shen2003} vs.\ \citealt{vanderWel2014}).
    \textit{Panel (b):} the ratio
    $\mathcal{R}(z)\!\equiv\!R_{\rm eff}(z)/[R_{0}(1+z)^{-3/4}]$,
    with HEAT's closed form (solid blue) asymptoting to
    $\Omega_{m}^{-1/4}\!\approx\!1.335$ (dotted) and the
    matter-era baseline $\mathcal{R}\!=\!1$ (dashed).  The orange
    marker in both panels is the unlensed three-disk subset of the
    ALMA kinematic sample at $z\!\approx\!4.5$
    ($N\!=\!3$, Roman-Oliveira+23/24; SPT0418-47 excluded as
    lensed).  The full four-disk sample is used for the kinematic
    cross-check in Fig.~\ref{fig:kin}; this three-object overlay is
    shown only as a supporting size-normalisation comparison rather
    than a precision constraint.  It lies near the HEAT curve and
    above the matter-era power law by a factor comparable to
    $\mathcal{R}_{\infty}$.  Current data are consistent with
    both curves within the $\sim 20$--$30\%$ systematic band
    and do not yet discriminate between them;
    Euclid/Roman/JWST mass-complete samples will reduce this
    band to the percent level, converting the figure into a
    quantitative test.\label{fig:norm}}
\end{figure}

Upcoming homogeneous, mass-selected size catalogues---Euclid
over $0.1\!\lesssim\!z\!\lesssim\!2$, Roman HLS to
$z\!\sim\!3$, and JWST mass-complete samples at $z\!\gtrsim\!3$---
are expected to reduce the per-bin uncertainty on
$\mathcal{R}(z)$ to the percent level, converting the present
consistency check into a quantitative test of the
$\azero\propto H(z)$ coupling and, given that coupling, of the
$\lcdm$ expansion history on structure scales, independently
of Type Ia supernovae and baryon acoustic oscillations.

\begin{figure}[t]
  \centering
  \includegraphics[width=\columnwidth]{figures/fig5_size_mass.pdf}
  \caption{Observed galaxy compactification versus redshift:
    $R_{\rm obs}/R_{0}(M_{\star},z\!=\!0)$ for spectroscopically
    confirmed high-$z$ galaxies, normalised to the local
    late-type size--stellar-mass relation
    \citep{vanderWel2014}.  Solid blue: scaling-only HEAT model
    curve $R/R_0=[H(z)/H_0]^{-1/2}$ (Eq.~\ref{eq:RofH}); the
    surrounding shaded band is the $\pm 20\%$ local-anchor
    systematic \citep{Shen2003,vanderWel2014}.  Green circles
    denote ALMA kinematic detections; purple diamonds are
    spectroscopic-$z$ systems with photometric sizes.  Vertical
    error bars propagate the tabulated $\sigma_{\log M_{\star}}$
    through $R_{0}\propto M_{\star}^{0.22}$; faded halos combine
    this in quadrature with a conservative $\pm 30\%$
    non-equilibrium systematic for galaxies flagged as
    merger, lensed, or dual in their data caveat, which is why
    SPT0418-47 and SGP38326 carry visibly larger uncertainty
    than CRISTAL-22 and DC-881725.  Orange squares: same-sample
    binned median $R_{e}/R_{0}$ from the MUSE-DARK MHUDF
    photometry catalogue (\citealt{Ciocan2026Halo} Paper~I,
    $\log_{10} M_{\star}/M_{\odot}\!>\!9$, three equal-count
    $z$-bins) with 16/84 percentile error bars on both axes.
    The intermediate-$z$ binned means from
    Table~\ref{tab:interz} and the MHUDF overlay are compared at
    the consistency level set by local-anchor systematics; points at
    $z\!\gtrsim\!10$ require non-equilibrium arguments (e.g.\
    incomplete virialisation, ongoing accretion) deferred to a
    longer treatment.\label{fig:size}}
\end{figure}

\subsection{MUSE-DARK RAR shape test at
  $z\approx 0.3$--$1.4$}\label{sec:a0z_priority}

Before using the absolute Milgrom--Hubble anchor for velocities
or BTFR zero-points, the direct MUSE-DARK RAR measurements provide
a cleaner test of the redshift \emph{shape}.  We therefore take
the one-parameter family $\azero(z)=K\,cH(z)$ as the main
comparison: it preserves all scaling-only predictions of
\S\ref{sec:exp}, including $R(z)/R_{0}$, while asking whether the
observed $\azero(z)$ points have the predicted $H(z)$ form.
On the clean $N\!=\!7$ cohort (\citealt{McGaugh2016},
\citealt{Desmond2023}, \citealt{Varasteanu2025} combined fit,
and the four MUSE-DARK Paper~III bins) this free-$K$ shape gives
$\chi^{2}\!=\!5.9$ for $6$ degrees of freedom with best-fit
$\beta\!\equiv\!K/K_{\rm HEAT}\!=\!1.29\!\pm\!0.04$ at Planck
cosmology.  The HEAT $n\!=\!1$ shape is accepted by the
likelihood-ratio test against the free-$n$ alternative at
$p\!=\!0.47$, and a $3\!+\!4$ predictive cross-validation under
HEAT $n\!=\!1$ is internally consistent in both directions
($\chi^{2}_{\rm fwd}\!=\!7.89$ vs.\ $\chi^{2}_{\rm crit}\!=\!9.49$;
$\chi^{2}_{\rm back}\!=\!3.97$ vs.\ $\chi^{2}_{\rm crit}\!=\!7.81$).
A uniform $+0.11$\,dex M/L revision applied to the four Ciocan
high-$z$ bins (within the Paper~III per-galaxy M/L uncertainty
floor) reduces the prefactor to $\beta\!=\!1.06\!\pm\!0.03$,
illustrating that the $\beta\!\neq\!1$ residual is jointly
degenerate with the stellar M/L calibration; the same revision
fails the forward predictive cross-validation, so the raw
cohort is the more self-consistent slice of the calibration
envelope.  The conventional $\azerolocal=cH_0/(2\pi)$ curve is
kept as the secondary absolute-reference line, not as the
definition of the headline size prediction.  The detailed
goodness-of-fit ladder, the six robustness tests
(MCMC posterior, cosmology grid, LRT, bootstrap, cross-validation,
M/L profiles), and the halo-density counterpart are given in
\S\ref{sec:a0z} and Figure~\ref{fig:a0z}.

\subsection{ALMA joint velocity--size consistency check at
  $z\sim 4.5$ ($N=4$, internal cross-check)}\label{sec:alma}

All four galaxies at $z>4$ with published resolved ALMA
rotation velocities---SPT0418-47 at $z=4.22$ \citep{Rizzo2020}
and three disks at $z\sim 4.5$
\citep{RomanOliveira2023,RomanOliveira2024}---admit an
\emph{internal consistency check} (not a statistical test, given
$N=4$) of the scaling in Eq.~\ref{eq:a0scale}.  Because
$v_{\rm flat}\propto\azero^{+1/4}$ and $R\propto\azero^{-1/2}$
scale with opposite powers of the acceleration, a joint match
on both axes probes the functional form of $\azero(z)$, not
merely its normalisation: a rescaling that improves one
observable generically worsens the other.  We treat the
resulting numbers as supporting consistency rather than as
independent evidence; the headline test remains the
$\sim 10^{4}$-galaxy size shape of \S\ref{sec:interz}.

Using published total baryonic masses and the conventional
absolute anchor of Eq.~\ref{eq:a0}, we compute the deep-MOND flat velocity
$v_{\rm flat}=(GM_b\,\azero)^{1/4}$ and the equilibrium
compactification $R/R_{0}=[\azero(z)/\azero(0)]^{-1/2}$ for
each galaxy (Table~\ref{tab:alma}).  The zero-parameter
sample-mean ratios are
$\langle v_{\heat}/v_{\rm obs}\rangle=1.00$ and
$\langle R_{\rm obs}/R_{0}\rangle_{\rm pred}=0.37$, jointly
consistent with the observed sample means ($1.00$ and $0.38$
respectively) on both axes within the $N{=}4$ sample-mean
envelope.  The corresponding constant-$\azero$ MOND
velocity mean is $0.63$.  Leave-one-out resampling gives
$1.00\pm 0.08$ (HEAT) versus $0.63\pm 0.05$ (MOND) on the
velocity axis; the HEAT advantage survives removal of any
single galaxy, although at $N\!=\!4$ jackknife uncertainty
estimates are themselves unreliable and should be read as
order-of-magnitude guides rather than formal confidence
intervals.  Per-galaxy $v_{\rm obs}$ values carry
intrinsic systematic uncertainties at the $15$--$20\%$ level
from inclination, beam smearing, and turbulent-broadening
subtraction, so the residual scatter in individual HEAT-to-
observed ratios ($0.75$--$1.36$) is consistent with a
zero-parameter prediction combined with a realistic
observational error budget.  Formal statistical significance
is precluded by $N=4$; the definitive test requires
$\gtrsim 10$ resolved rotation curves at $z>3$, which
upcoming ALMA large programmes will provide.

\begin{deluxetable}{lcccc}
\tablecaption{Joint zero-parameter HEAT kinematic and size
  predictions for the four ALMA disks at
  $z\!\sim\!4.5$.\label{tab:alma}}
\tablehead{
  \colhead{Galaxy} & \colhead{$z$} &
  \colhead{$v_{\heat}/v_{\rm obs}$} &
  \colhead{$R_{\rm obs}/R_{0}$} &
  \colhead{$R_{0}$-pred}
}
\startdata
CRISTAL-22       & 4.53 & 1.00 & 0.37 & 0.37 \\
DC-881725        & 4.56 & 1.36 & 0.32 & 0.37 \\
SGP38326         & 4.42 & 0.75 & 0.45 & 0.37 \\
SPT0418-47       & 4.22 & 0.89 & 0.18 & 0.38 \\
\hline
Mean (all 4)     & ---  & 1.00 & ---  & ---  \\
Mean (excl.\ SPT)& ---  & ---  & 0.38 & 0.37 \\
\enddata
\tablecomments{$v_{\heat}/v_{\rm obs}$: ratio of HEAT deep-MOND
  $v_{\rm flat}=(GM_b\azero)^{1/4}$ to observed ALMA velocity,
  using the conventional absolute anchor
  $\azero(0)=cH_0/(2\pi)$.
  $R_{\rm obs}/R_{0}$: observed size normalised to the local
  late-type size--stellar-mass relation
  \citep{vanderWel2014}.  $R_{0}$-pred:
  $[\azero(z)/\azero(0)]^{-1/2}$.  SPT0418-47 is a lensed
  ongoing minor merger with companion SPT0418B at $4.42$\,kpc
  projected separation \citep{Cathey2024}; the size is the
  \mbox{de-lensed} intrinsic [C\,II] half-light radius of
  \citet{Rizzo2020} and the system is excluded from the
  size-ratio mean because disk-equilibrium scaling is not
  expected to hold during an interaction.  The kinematic axis
  is retained because the $V/\sigma=9.7$ disk rotation measured
  by \citet{Rizzo2020} is robust to the merger.}
\end{deluxetable}

\begin{figure}[t]
  \centering
  \includegraphics[width=\columnwidth]{figures/fig3c_kinematic.pdf}
  \caption{Sample-mean ALMA kinematic cross-check ($N{=}4$).
    Observed rotation
    velocities (points) versus locally anchored HEAT and
    constant-$\azero$ MOND model predictions (curves) for the four
    $z>4$ disks in Table~\ref{tab:alma}.  The HEAT curve is
    consistent with the sample mean within the stated small-sample
    and observational-systematic envelope; the MOND curve falls
    below it by a factor $1/0.63\!\approx\!1.6$.  Given
    $N{=}4$ this is an internal consistency check rather than a
    statistical test.
    \label{fig:kin}}
\end{figure}

\subsection{BTFR zero-point: consistency with the predicted
  trend}\label{sec:btfr}

Combining the deep-MOND baryonic Tully--Fisher relation
$V_{\rm flat}^{4}=GM_{b}\,\azero$
\citep{Milgrom1983,FamaeyMcGaugh2012} with the
scaling ansatz Eq.~\ref{eq:a0scale} implies a locally anchored
shift of
the BTFR zero-point at fixed rotation velocity,
\begin{equation}
  \Delta\log_{10}M_{b}(z)=-\log_{10}\!\left[\frac{H(z)}{H_{0}}\right],
  \label{eq:btfr-evo}
\end{equation}
while constant-$\azero$ MOND predicts
$\Delta\log_{10}M_{b}(z)=0$ independently of redshift.  As with
the size ratio $\mathcal{R}(z)$ of \S\ref{sec:exp}, we use the
presently available heterogeneous sample only to check whether
the direction and amplitude of the predicted trend are
consistent with the data; a decisive quantitative test awaits a
homogeneous cosmic-noon rotation-curve sample.  Projecting the
four $z>4$ ALMA kinematic sources of Table~\ref{tab:alma} onto
the axis Eq.~\ref{eq:btfr-evo} (Figure~\ref{fig:btfr}) yields
$\Delta_{\rm obs}$ values of $-0.87$ dex (CRISTAL-22, $z=4.53$),
$-0.33$ dex (DC-881725, $z=4.56$), and $-1.03$ and $-1.36$ dex
for the caveated SPT0418-47 and SGP38326 systems, against a
HEAT prediction of $\simeq-0.86$ dex at $z\!\sim\!4.5$.  The SPARC
$z=0$ anchor (171 galaxies, \citealt{Lelli2016}) sits on the
HEAT curve within the \citet{Lelli2019} intrinsic scatter of
$\sim 0.10$ dex, by construction of the local BTFR anchor;
CRISTAL-22 --- the most rotation-dominated
equilibrium system in the high-$z$ sample --- sits on the
scaling prediction to within the per-galaxy
systematics.  The remaining points scatter above and below
the curve at amplitudes consistent with the known
merger/lensed/clumpy-disk systematics, and the overall trend
is in the direction predicted by Eq.~\ref{eq:btfr-evo} and
opposite to the constant-$\azero$ MOND baseline.  A homogeneous
cosmic-noon ($z\sim 1$--$2.5$) sample of $\sim 10$
rotation-dominated disks with individual $\lesssim 15\%$
velocity precision would convert this trend consistency into a
quantitative discrimination between HEAT, CEG, and
constant-$\azero$ MOND at $>3\sigma$ at fixed baryonic mass.

\paragraph{Tension with the lensed bTFR zero-evolution result at
  $z\!\sim\!1$.}
\citet{Jeanneau2026} (MUSE-DARK Paper II) report no detectable
evolution of the lensed bTFR zero-point at $z\!\sim\!1$,
$\Delta\log_{10}M_{b}\!=\!0.00\!\pm\!0.05$\,dex, against the
HEAT prediction $-\log_{10}[H(z\!=\!1)/H_{0}]\!\approx\!-0.26$\,dex.
Taken at face value this is a $\sim 4\sigma$ challenge to
Eq.~\ref{eq:btfr-evo} and we list it as such in
\S\ref{sec:falsify}.  Several systematics temper the immediate
interpretation but do not dissolve the tension; we adopt here
the interpretation outlined to us by the MUSE-DARK collaboration
(B.~Ciocan, priv.\ comm., 2026) together with the internal
discussions in Papers~II and~III.  The dominant uncertainty on
\emph{both} sides of the comparison is the baryonic-mass budget,
but the leading component differs.  In Paper~II the lensed bTFR
baryonic masses are reconstructed via scaling relations whose
calibrations carry non-negligible systematics, in particular the
HI mass and the \citet{Tacconi2018}-type molecular-gas scaling
relation that is extrapolated below its
well-calibrated stellar-mass range in Paper~II's regime; a coherent systematic shift
in the inferred baryon budget could plausibly drive an apparent
bTFR zero-point non-evolution, and independent cosmic-noon
measurements (e.g.\ \citealt{Ubler2017}) do not converge on a
single sign for that shift.  In addition, the Jeanneau
zero-point is derived assuming a \emph{static} BTFR slope and
$z\!=\!0$ scatter, which absorbs part of any redshift dependence
into the inferred normalisation, and individual lensed
baryonic-mass estimates carry a $\sim 0.2$\,dex source-plane
and mass-to-light systematic of the kind discussed in the
MUSE-DARK Paper~III appendix material, leaving room for a
coherent zero-point bias whose sign and magnitude cannot be
fixed from the present data alone.

In Papers~I and~III, the disk--halo decomposition uses parametric
HI models with priors on $M_{\rm HI}$ from comparable scaling
relations and---like the SPARC analyses---does not model molecular
gas as a separate component; the molecular contribution is
absorbed into the disk, contributing an additional systematic of
order $0.15$--$0.20$\,dex to the disk masses.  Paper~III shows
internally that a perturbation of this size on its $\azero(z)$
inversion is insufficient to reconcile the data with a
non-evolving $\azero$, so the Paper~II/Paper~III tension cannot
be removed by uniformly inflating the disk-mass error budget on
the supportive side alone.

Where we do not invoke these systematics, the direct $\azero(z)$
measurement of \citet{Ciocan2026RAR} (Paper~III), discussed in
\S\ref{sec:a0z}, is consistent with the HEAT \emph{shape}
$\azero(z)\!\propto\!H(z)$: on the clean $N\!=\!7$ cohort the
free-$K$ HEAT fit gives $\chi^{2}\!=\!5.9$ for $6$ dof, the
likelihood-ratio test for $n\!=\!1$ vs.\ free $n$ accepts
$n\!=\!1$ at $p\!=\!0.47$, and the $3\!+\!4$ predictive
cross-validation is internally consistent in both directions
at $95\%$.  The Ciocan MOND-framework $z\!=\!0$ intercept also
agrees with HEAT's zero-point (Paper~III App.~E) to within
$1\%$.  The trilogy as a whole therefore supports the
$\azero\!\propto\!H(z)$ coupling on at least two of three
independent observables; the Jeanneau bTFR zero-evolution result
remains the single most demanding falsifier of HEAT in the
current literature.  The asymmetry in our treatment of the supportive
vs.\ challenging MUSE-DARK observables is
observable-driven: Papers~I and~III derive $\azero(z)$ from
\emph{rotation-curve} modelling, where the kinematic observable
is independent of the lens model, whereas Paper~II's bTFR
zero-point depends directly on the lensed baryonic-mass
reconstruction.  A direct HI-based bTFR measurement at
$z\!\sim\!1$, ultimately requiring SKA-precursor and very deep
ALMA programmes that resolve cold gas without extrapolation
(\S\ref{sec:falsify}, test~7), would resolve this
observable-level distinction.

\begin{figure}[t]
  \centering
  \includegraphics[width=\columnwidth]{figures/fig7_btfr_evolution.pdf}
  \caption{BTFR zero-point offset
    $\Delta\log_{10}M_{b}$ at fixed $V_{\rm flat}$ versus
    redshift.  Solid blue: locally anchored HEAT model prediction
    $-\log_{10}[H(z)/H_{0}]$ (Eq.~\ref{eq:btfr-evo}); surrounding
    shaded band is the $\pm 0.10$ dex intrinsic BTFR scatter
    from \citet{Lelli2019}.  Dashed grey: constant-$\azero$ MOND
    baseline.  Black square: SPARC $z=0$ median with 16/84
    percentile uncertainty ($N{=}171$, \citealt{Lelli2016}).
    Orange triangles: cosmic-noon binned BTFR offsets from
    \citet{Ubler2017} at $z\!\sim\!0.9$ and $z\!\sim\!2.3$.
    Green circles: equilibrium ALMA disks at $z\!\sim\!4.5$;
    purple diamonds: merger/lensed/dual systems retained for
    completeness.  The red cross marks the
    \citet{Jeanneau2026} (MUSE-DARK Paper II) zero-evolution
    measurement at $z\!\sim\!1$ from a strongly lensed bTFR sample
    ($\Delta\log_{10}M_{b}\!=\!0.00\!\pm\!0.05$\,dex), in
    tension with the HEAT-predicted shift
    $-\log_{10}[H(z\!=\!1)/H_{0}]\!\approx\!-0.26$\,dex.  Error
    bars on the ALMA points combine the SED-based
    $\sigma_{\log M_{\star}}$ with a $10\%$ systematic on
    $V_{\rm flat}$ (i.e.\ $\sim 0.17$ dex on $\log V^{4}$).
    The heterogeneous cosmic-noon and high-$z$ ALMA points are used
    only as a trend-level comparison; the Jeanneau+2026
    zero-evolution result is discussed in \S\ref{sec:btfr} and
    listed as a falsification target in \S\ref{sec:falsify}.
    \label{fig:btfr}}
\end{figure}

\subsection{MUSE-DARK RAR: detailed $\azero(z)$
  shape and anchor tests at intermediate
  redshift}\label{sec:a0z}

The redshift evolution of $\azero$ is the single feature on
which HEAT differs from constant-$\azero$ MOND.  The MUSE-DARK
collaboration has now provided a direct measurement of $\azero$
at $0.33\!<\!z\!<\!1.44$ from intermediate-redshift radial
acceleration relation (RAR) fits to a sample of MUSE
spectro-imaging targets in the Hubble Ultra-Deep Field
\citep{Ciocan2026RAR}.  We compare the resulting $\azero(z)$
points against four reference scalings: (i) the free-$K$ HEAT
shape $\azero(z)=K\,cH(z)$, which tests whether the data follow
the predicted $H(z)$ form independently of the absolute anchor;
(ii) the conventional HEAT absolute anchor
$\azero(z)\!=\!\azerolocal H(z)/H_0$ with
$\azerolocal=cH_0/(2\pi)$; (iii) the constant
$\azero(0)\!\equiv\!1.20\!\times\!10^{-10}\,\mathrm{m\,s^{-2}}$
of \citet{McGaugh2016}; and (iv) the linear-in-$z$
parameterisations $\azero(z)\!=\!a_{0}(0)\!+\!a_{1}\,z$ that
Paper~III adopts as a phenomenological summary in three
independent fitting frameworks: a global DC14 prescription
(adopted because DC14 emerges as the best-fitting profile across
the Paper~I sample), a per-galaxy preferred-profile fit in which
each galaxy is assigned the halo family selected by pairwise
Bayesian model comparison from the six profiles tested in
Paper~I (NFW, cNFW, Burkert, Dekel--Zhao, DC14, and Einasto), and
a direct MOND-framework fit.

\paragraph{Clean cohort.}  The $z\!=\!0$ anchor combines the
SPARC RAR-fit median of \citet{McGaugh2016}
($1.20\!\pm\!0.26\!\times\!10^{-10}\,\mathrm{m\,s^{-2}}$) with
the tighter full-Hamiltonian Monte Carlo joint inference of
\citet{Desmond2023} on the same SPARC sample
($1.19\!\pm\!0.10$, MLS interpolating function, matching the
framework adopted by \citealt{Ciocan2026RAR}); we use the
\citet{Varasteanu2025} combined MIGHTEE-HI + high-acceleration
SPARC RAR-framework value
($1.32\!\pm\!0.13$ at $z\!\sim\!0.04$) as the lowest
intermediate-$z$ anchor.  Together with the four
\citet{Ciocan2026RAR} bins at $z\!\sim\!0.55$, $0.78$, $1.05$,
$1.30$ this defines the clean $N\!=\!7$ cohort on which we report
all goodness-of-fit numbers in this section.  Numerical outputs
are reproducible from
\texttt{publication/a0\_ml\_degeneracy.py} (Sec.~4.5 fitting
ladder) and \texttt{publication/a0\_robustness\_tests.py}
(robustness tests A--F summarised below).

\paragraph{Headline goodness of fit on the raw cohort.}  Allowing
only the proportionality constant $K$ in
$\azero(z)\!=\!K\,c\,H(z)$ to float yields a best-fit
$\beta\!\equiv\!K/K_{\rm HEAT}\!=\!1.29\!\pm\!0.04$ with
$\chi^{2}\!=\!5.9$ for $6$ degrees of freedom (Planck-2018
$H_{0}\!=\!67.4\,\mathrm{km\,s^{-1}\,Mpc^{-1}}$, $\Omega_{m}\!=\!0.315$).
The constant-$\azero$ MOND baseline gives $\chi^{2}\!=\!198$ for
the same cohort and is excluded at $\Delta\chi^{2}\!>\!190$
against the free-$K$ HEAT shape.  The three Ciocan linear-in-$z$
fits give $\chi^{2}\!=\!13$--$18$ on the same cohort, worse than
the free-$K$ HEAT shape despite carrying one extra fitted
parameter; the full goodness-of-fit ladder is collected in
Table~\ref{tab:a0z_chi2}.

\paragraph{Test of the $\azero(z)\!\propto\!H(z)$ shape.}  The
specific HEAT prediction is not the absolute prefactor $\beta$
(which is jointly degenerate with $H_{0}$ and the stellar
mass-to-light calibration, as quantified below), but the shape
exponent: $\azero(z)\!\propto\!H(z)^{n}$ with $n\!=\!1$.  We test
this by promoting $n$ to a free parameter and applying the
likelihood-ratio test (LRT) for the nested $n\!=\!1$ hypothesis.
On the raw cohort the free-$n$ fit recovers
$n\!=\!+1.08\!\pm\!0.10$ at $\chi^{2}\!=\!5.42$ for $5$ degrees
of freedom; the $n\!=\!1$ vs.\ $n$-free likelihood ratio gives
$\Delta\chi^{2}\!=\!0.52$ for one extra parameter
($p\!=\!0.47$, $0.7\sigma$).  The data therefore do not reject
the HEAT $n\!=\!1$ shape.  Independently, the MCMC marginal
posterior on $n$ (Test A of
\texttt{a0\_robustness\_tests.py}) gives
$n\!=\!1.07\!\pm\!0.11$ ($68\%$ credible region containing
$n\!=\!1$), and a $10^{4}$-draw bootstrap (Test D) returns the
$95\%$ non-parametric interval $n\!\in\![0.78, 1.44]$
(containing $n\!=\!1$).  The Hubble-scaling agreement is also
visible directly in the data: anchored at the
\citet{Desmond2023} $z\!=\!0$ value, the ratio
$[\azero^{\rm obs}(z)/\azero^{\rm SPARC}]/[H(z)/H_{0}]$ stays
within $1.22$--$1.40$ across the four \citet{Ciocan2026RAR}
bins, with no significant residual trend in $z$.

\paragraph{Predictive cross-validation.}  The LRT and MCMC tests
above are likelihood-based and do not by themselves test
\emph{predictive} consistency.  We therefore perform a
$3\,{+}\,4$ split of the cohort and evaluate the HEAT
$n\!=\!1$ extrapolation in both directions (Test E of
\texttt{a0\_robustness\_tests.py}).
\textit{Forward direction}: training only on the three $z\!\lesssim\!0.04$
anchors gives $\beta\!=\!1.18\!\pm\!0.07$; the four held-out
\citet{Ciocan2026RAR} bins at $z\!\sim\!0.55$--$1.30$, evaluated
against this low-$z$-only HEAT model, give a predictive
$\chi^{2}\!=\!7.89$ for $4$ degrees of freedom against the
$95\%$-critical value $\chi^{2}_{\rm crit}\!=\!9.49$
(accepted at $95\%$).  \textit{Back direction}: training only on
the four Ciocan high-$z$ bins (with $n$ fixed at the HEAT value
$n\!=\!1$) gives $\beta\!=\!1.33\!\pm\!0.05$; the three held-out
low-$z$ anchors, evaluated against this high-$z$-only HEAT
model, give $\chi^{2}\!=\!3.97$ for $3$ degrees of freedom
against $\chi^{2}_{\rm crit}\!=\!7.81$ (also accepted at
$95\%$).  The cohort is bidirectionally consistent with the
HEAT shape: the low-$z$ anchor predicts the MUSE-DARK trend,
and the MUSE-DARK trend back-extrapolates to the low-$z$ anchor,
under a single common $\beta$ in each direction whose internal
$1\sigma$ values ($1.18$, $1.33$) bracket the joint fit
($1.29$).  (Allowing $n$ to float on the Ciocan-only subset
inflates the back-direction $\chi^{2}$ to $29$; this is a
two-parameter-on-four-points degeneracy artifact -- the
$n\!=\!1$ fit on the same four Ciocan bins is only
$\Delta\chi^{2}\!=\!2.3$ worse than the $n$-free best fit --
not a HEAT failure.)

\paragraph{Common-framework $\chi^{2}$.}  All three low-$z$
anchors and the four MUSE-DARK bins share the same underlying
methodology: a standard $\Lambda$CDM disk--halo decomposition of
the rotation curve, supplying $a_{\rm tot}$ from the HI velocity
field and $a_{\rm bar}$ from the baryonic surface-density
profile via Poisson's equation, with a MOND-flavoured
interpolation function fit \emph{only} to the resulting
$a_{\rm tot}$ vs.\ $a_{\rm bar}$ relation
\citep[as explicitly stated for SPARC by][and adopted
identically in Paper~III]{McGaugh2016}.  The
\citet{Varasteanu2025} combined fit reports $\azero$ from this
same RAR framework applied jointly to their MIGHTEE-HI sample
and the high-acceleration portion of SPARC, not from a primary
bTFR inversion.  We therefore treat all seven points as
methodology-consistent and do not partition the cohort by
inversion-type.  The residual methodological lever remaining is
which halo-density prescription is allowed at the disk--halo
decomposition step inside Paper~III itself: a global DC14
template, per-galaxy preferred profiles, or a direct
MOND-framework fit.  These choices shift the recovered
$\azero(z\!=\!1)$ across $2.19$--$2.61\!\times\!10^{-10}\,
\mathrm{m\,s^{-2}}$ (Eqs.~2, D.1, E.4 of Paper~III) and define
the $\chi^{2}$ spread among the three Ciocan linear-in-$z$
parameterisations in Table~\ref{tab:a0z_chi2}.

\begin{deluxetable*}{lccc}
\tablecaption{Goodness-of-fit comparison for $\azero(z)$
  parameterisations on the clean $N\!=\!7$ cohort
  (\citealt{McGaugh2016} and \citealt{Desmond2023} at
  $z\!=\!0$, the combined MIGHTEE-HI\,+\,high-$g$-SPARC RAR-fit
  of \citealt{Varasteanu2025} at $z\!\sim\!0.04$, and the four
  \citealt{Ciocan2026RAR} bins at
  $z\!\sim\!0.55$--$1.30$), reported both raw and after the
  uniform $+0.11$\,dex M/L correction discussed below
  (\S\ref{sec:a0z}, ``Stellar-mass systematic'').  The free-$K$
  HEAT shape gives the lowest $\chi^{2}$ among the entries
  shown, including each of the three Ciocan linear-in-$z$
  parameterisations \emph{despite carrying one fewer fitted
  parameter} (1 vs.\ 2).  The constant-$\azero$ MOND baseline is
  strongly disfavoured at $\chi^{2}\!\gtrsim\!95$ on both
  columns.  Note that under the $+0.11$\,dex M/L correction the
  free-$K$ fit recovers $\beta\!=\!1.06\!\pm\!0.03$ (the bare
  HEAT prefactor at $\sim 2\sigma$) but the same correction
  fails the forward $3\!\to\!4$ predictive cross-validation on
  the raw cohort (\S\ref{sec:a0z}, ``Stellar-mass systematic''),
  so the M/L column is a calibration sensitivity rather than a
  preferred analysis.\label{tab:a0z_chi2}}
\tablewidth{0pt}
\tablehead{
  \colhead{Model} & \colhead{Fitted pars} &
  \colhead{$\chi^{2}_{\rm raw}$ ($N\!=\!7$)} &
  \colhead{$\chi^{2}_{+0.11\,\rm dex\,M/L}$ ($N\!=\!7$)}
}
\startdata
HEAT shape, free $K$                & 1    & {\bf 5.9}  & {\bf 6.2}  \\
HEAT, $\azerolocal\!=\!cH_0/(2\pi)$ & 0    & 61.8       & 9.7  \\
Constant-$\azero$ MOND ($1.20$)     & 0    & 198.4      & 95.1 \\
Ciocan DC14 uniform                 & ext. & 13.0       & $104.7^{\dagger}$ \\
Ciocan best-fit per-galaxy           & ext. & 12.5       & $129.9^{\dagger}$ \\
Ciocan MOND framework               & ext. & 17.6       & $29.1^{\dagger}$  \\
\enddata
\tablecomments{``Fitted pars'' is the number of parameters
  refit on the $N\!=\!7$ cohort of this work; the free-$K$ fit
  therefore has $6$ dof.  Rows marked ``ext.'' are external
  two-parameter linear fits to the broader
  \citet{Ciocan2026RAR} sample (Eqs.~4, D.2, E.5 of
  Paper~III); their slopes and intercepts are fixed for our
  $\chi^{2}$ accounting and no dof reduction is applied.
  ``HEAT shape, free $K$'' refits only the single proportionality
  constant in $\azero(z)=K\,cH(z)$, finding $\beta\!\equiv\!
  K/K_{\rm HEAT}\!=\!1.29\!\pm\!0.04$ raw and $1.06\!\pm\!0.03$
  after the $+0.11$\,dex M/L correction; the headline observable
  $R(z)/R_{0}$ is invariant under this prefactor, \S\ref{sec:exp}.
  ``HEAT, $\azerolocal\!=\!cH_0/(2\pi)$'' uses the conventional
  absolute anchor of Eq.~\ref{eq:a0} ($\beta\!=\!1$) and is
  retained as the secondary normalisation reference; its
  $\chi^{2}_{\rm raw}\!=\!62$ falls to $9.7$ under the same
  M/L correction, i.e.\ becomes statistically indistinguishable
  from the free-$K$ fit at $\sim 2\sigma$, approaching the bare
  HEAT prediction.
  $^{\dagger}$ Entries marked with a dagger are the
  \emph{unrefitted} Ciocan linear $\azero(z)\!=\!a_0\!+\!a_1\,z$
  parameterisations evaluated against the M/L-corrected
  data; the corrected data want a markedly shallower slope, so
  the unrefitted DC14 and per-galaxy linear fits fail
  $\chi^{2}\!\gtrsim\!100$ in the second column.  These dagger
  entries should be read as ``this $a_1$ slope is incompatible
  with the M/L-corrected data,'' not as competitive model fits.
  The Ciocan MOND-framework intercept
  $\azero^{\rm MOND}(0)\!=\!1.03\!\pm\!0.05$ remains the
  apples-to-apples $z\!=\!0$ anchor for any cohort whose high-$z$
  points are themselves derived in the Ciocan-MOND framework.
  Numbers are reproduced from \texttt{a0\_evolution\_stats.txt}
  and \texttt{a0\_ml\_degeneracy.txt} in the Zenodo deposit and
  visualised in
  Figure~\ref{fig:a0z}.}
\end{deluxetable*}

\paragraph{Cross-framework normalisation match.}  In the MOND
framework adopted in Paper~III Appendix~E, the linear-fit
intercept is $\azero(z\!=\!0)\!=\!1.03\!\pm\!0.05\!\times\!10^{-10}\,
\mathrm{m\,s^{-2}}$.  The Milgrom--Hubble anchor adopted for
HEAT's absolute tests at
$z\!=\!0$ is $cH_{0}/(2\pi)\!=\!1.04\!\times\!10^{-10}\,
\mathrm{m\,s^{-2}}$.  These two numbers agree to within $1\%$,
and this is the apples-to-apples comparison: both invert the
same MOND interpolation function on the same observed
kinematics, so any methodology-dependent normalisation drift
is differenced out.  The DC14 uniform fit
($1.00\!\pm\!0.04$) and the best-fit per-galaxy fit
($1.05\!\pm\!0.05$) bracket the HEAT value within their
$1\sigma$ envelope, again with the MOND interpolation
function used inside the disk--halo decomposition.

The SPARC RAR-fit value
$\azero^{\rm SPARC}(0)\!=\!1.20\!\pm\!0.26$ is a different
inversion of partially overlapping physics: a galaxy-by-galaxy
RAR median rather than a population-level interpolation-function
fit.  The exchange between \citet{Rodrigues2018} and
\citet{McGaugh2018} documents methodology-dependent shifts of
order $\sim 30\%$ between the two reductions on the same SPARC
subsample, of the same
direction and roughly the same magnitude as the
$\sim 13\%$ offset between HEAT's $1.04$ and SPARC's $1.20$.
The adopted anchor therefore matches the framework-self-consistent
(MOND-interpolation) $z\!=\!0$ anchor at the level
quantified above and remains within $1\sigma$ of the SPARC RAR-fit
anchor; we treat the two anchors as bounding the
methodology-dependent normalisation rather than as a single
common datum.

\paragraph{Appendix robustness.}  The positive $\azero(z)$
trend in Paper~III is not tied to a single halo
parameterisation.  Paper~III's appendix repeats the inversion
under a uniform DC14 prescription, best-fit per-galaxy fits, and a
MOND-framework fit; all three recover an increasing acceleration
scale across the MUSE-DARK redshift interval, with a common
qualitative shape but methodology-dependent intercepts.  This
is why Fig.~\ref{fig:a0z} carries those reductions as external
calibration curves rather than as independent data points, and
why Table~\ref{tab:a0z_chi2} reports their $\chi^{2}$ alongside
both the free-$K$ HEAT shape test and the conventional-anchor
HEAT reference.

\paragraph{Stellar-mass systematic.}  The natural objection
\citep[discussed by][see also Paper~III \S\,4]{Ciocan2026RAR}
is in the opposite direction from a naive expectation: the
MUSE-DARK dynamical disk--halo decompositions are potentially
\emph{submaximal}, so the high-$z$ baryonic mass
$M_{\star}\!+\!M_{\rm HI}$ may be \emph{under-estimated}, which
in the deep-MOND limit
$g_{\rm obs}\!=\!\sqrt{g_{\rm bar}\azero}$ biases the recovered
$\azero$ \emph{upward} by a factor
$10^{+\delta_{M_{\star}}}$.  Paper~III Appendix~C tests this
along the trend axis by applying coherent $+0.2$ and $+0.45$\,dex
shifts to the disk mass-to-light ratio of the entire sample;
the recovered positive $\azero(z)$ \emph{slope} persists in both
cases, and the Paper~I appendix further demonstrates the
MUSE-DARK dynamical $M_{\star}$ are consistent with
SED-derived stellar masses at the same redshift.  Paper~III's
\emph{trend} conclusion therefore stands independent of the
M/L assumption.

The complementary question we ask here is at what
$\delta_{M_{\star}}$ the empirical \emph{prefactor}
$\beta\!\equiv\!K/K_{\rm HEAT}$ relaxes to unity.  Applying a
uniform $\delta_{M_{\star}}$ correction only to the four
Ciocan high-$z$ bins (the SPARC and Varasteanu+2025 anchors use
independently calibrated M/L), we find the high-$z$ weighted
$K$ reaches unity at
$\delta_{M_{\star}}\!=\!+0.124\!\pm\!0.05$\,dex and the
low-$z$/high-$z$ aggregate $K$ values reconcile at
$\delta_{M_{\star}}\!=\!+0.05$\,dex; both lie within the
per-galaxy M/L uncertainty floor quoted in Paper~III and below
the $+0.2$ to $+0.45$\,dex range that Paper~III Appendix~C
tested for the trend.  Refitting the cohort under a fiducial
$+0.11$\,dex uniform correction yields
$\beta\!=\!1.06\!\pm\!0.03$ and reduces the
$\beta\!=\!1$ residual from $\chi^{2}\!=\!62$ to
$\chi^{2}\!=\!9.7$.

We note, however, that the $+0.11$\,dex uniform correction is
not preferred by the predictive cross-validation of the
preceding paragraph: applying the same correction to the four
Ciocan bins and re-running the forward direction
(3 low-$z$ anchors $\to$ 4 high-$z$ bins under HEAT $n\!=\!1$)
gives predictive $\chi^{2}\!=\!10.14$ for $4$ degrees of
freedom, just above the $95\%$ critical value of $9.49$
(rejected at $95\%$; Test E2 of
\texttt{a0\_robustness\_tests.py}).  The uniform correction
shifts the four Ciocan bins below the band predicted by the
low-$z$ anchor; the residuals all turn negative
($-0.89, -1.15, -1.83, -2.16\sigma$ at the four bins), in
contrast to the raw-cohort positive residuals
($+1.91, +1.85, +0.82, +0.37\sigma$).  We therefore treat the
$+0.11$\,dex correction as a sensitivity test for the
prefactor, not as a preferred analysis: the raw cohort with
the bare HEAT shape is the more self-consistent fit when the
goal is bidirectional predictivity.  The
$(\beta, n)$ region tabulated below is the
calibration-marginalised envelope of plausible analyses; the
joint Planck-cosmology fit prefers the raw end of that
envelope.

\paragraph{$(\beta,n)$ calibration sensitivity.}  Repeating the
two-parameter fit
$\azero(z)\!=\!\beta\!\cdot\!cH_0/(2\pi)\!\cdot\![H(z)/H_0]^{n}$
under five well-motivated calibration choices traces a narrow
one-parameter ridge in $(\beta,n)$ space (Table~\ref{tab:beta_n_degeneracy}).

\begin{deluxetable}{lcc}
\tablecaption{$(\beta,n)$ calibration sensitivity on the clean
  $N\!=\!7$ cohort.  At fixed M/L the joint posterior is much
  tighter than the spread shown here (Test A of
  \texttt{a0\_robustness\_tests.py}: marginal posteriors give
  $\sigma_{\beta}\!\sim\!0.06$ and $\sigma_{n}\!\sim\!0.11$);
  what this table reports is the shift of the joint best fit
  under reasonable calibration variations.  All five rows
  exclude the constant-$\azero$ baseline ($n\!=\!0$) at
  $\Delta\chi^{2}\!\geq\!90$; none in isolation excludes
  $\beta\!=\!1$.\label{tab:beta_n_degeneracy}}
\tablehead{
  \colhead{Calibration} & \colhead{$\beta$} & \colhead{$n$}
}
\startdata
default (Ciocan bins, raw)            & $1.29\!\pm\!0.04$ & $+1.08\!\pm\!0.10$ \\
Ciocan MOND-IF global (raw)           & $1.18\!\pm\!0.05$ & $+1.00\!\pm\!0.13$ \\
$+0.06$\,dex M/L                      & $1.16\!\pm\!0.03$ & $+0.88\!\pm\!0.10$ \\
$+0.11$\,dex M/L                      & $1.06\!\pm\!0.03$ & $+0.72\!\pm\!0.10$ \\
$+0.11$\,dex M/L $+$ MOND-IF          & $0.99\!\pm\!0.04$ & $+0.56\!\pm\!0.13$ \\
\enddata
\tablecomments{The MOND-IF column uses the single Paper~III
  Eq.~E.4 sample-wide intercept ($\azero(z\!=\!1)\!=\!2.19$)
  in place of the four binned bin1--bin4 points, because that
  reduction inverts the same MOND interpolation function as the
  HEAT prediction itself.  The $+0.11$\,dex M/L column applies a
  uniform stellar mass-to-light revision only to the four Ciocan
  high-$z$ bins (and to the MOND-IF global in the last row),
  leaving the SPARC and Varasteanu+2025 anchors unchanged.
  The constant-$\azero$ baseline (the $n\!=\!0$ limit) is
  excluded at $\Delta\chi^{2}\!\geq\!90$ on every row.  Numbers
  are reproduced from \texttt{a0\_ml\_degeneracy.txt}.}
\end{deluxetable}

\noindent The five rows trace a calibration-driven
one-parameter family in $(\beta,n)$ space spanning
$(1.29,1.08)\!\to\!(0.99,0.56)$.  At fixed calibration the
joint posterior is narrow (Test A), but
the choice of M/L calibration shifts the joint best fit by
$\Delta\beta\!\sim\!0.3$ and $\Delta n\!\sim\!0.5$.  Three
independent statements survive this calibration sensitivity:
(i) the constant-$\azero$ alternative ($n\!=\!0$) is excluded
by every row at $\Delta\chi^{2}\!\geq\!90$ -- this is the
robust constant-vs.-evolving statement;
(ii) on the raw cohort the joint best fit is
$(\beta, n)\!=\!(1.29, 1.08)$, with the HEAT $n\!=\!1$ shape
accepted by the LRT at $p\!=\!0.47$ and by the bidirectional
cross-validation at $95\%$ in both directions;
(iii) the $\beta\!=\!1$ row of the table is reached only when
one of the M/L corrections is applied, and that same correction
fails the forward predictive cross-validation at $95\%$ on the
raw cohort -- so $\beta\!=\!1$ is accessible within the
calibration envelope but is not the most predictively
consistent slice of it.  We therefore quote the raw
$(\beta, n)\!=\!(1.29, 1.08)$ as the headline joint fit at
Planck cosmology, and the $(\beta\!=\!1, n\!\sim\!0.6$--$0.7)$
end as the M/L-corrected sensitivity.  Independent breaking of
this residual $(\beta,n)$ degeneracy requires either an
independent stellar M/L calibration at $z\!>\!0.5$
(e.g.\ JWST NIRSpec stellar populations on a
MUSE-DARK-overlap sample) or new $\azero(z)$ measurements at
the currently-empty $0.1\!<\!z\!<\!0.5$ range.  An $H_{0}$
sensitivity test on a $9{\times}5$ grid in
$H_{0}\!\in\![60,76]$ and $\Omega_{m}\!\in\![0.27,0.35]$
(Test B of \texttt{a0\_robustness\_tests.py}) further shows
that the joint $\chi^{2}$ for the bare HEAT
($\beta\!=\!1, n\!=\!1$) prediction falls from $\chi^{2}\!=\!62$
at Planck to $\chi^{2}\!=\!36$ at $H_{0}\!=\!73$; combined with
the $+0.11$\,dex M/L lever it falls to $\chi^{2}\!=\!8$ at
$H_{0}\!=\!73$.  In other words, $\beta\!=\!1$ becomes the
data-preferred slice only when both the SH$_{0}$ES-style $H_{0}$
and the M/L revision are taken together; neither lever alone is
sufficient.

\paragraph{Halo-density counterpart from Paper~I.}  An
independent corroboration comes from the same MUSE-DARK sample
analysed in the DM frame.  \citet{Ciocan2026Halo} (Paper~I) fit
six DM density profiles to the full sample of resolved
rotation curves at $z\!\sim\!0.85$ and report
$\rho_{s}\!\propto\!(1\!+\!z)^{0.54\pm 0.31}$ for the NFW-equivalent
characteristic density, i.e.\ a $\Delta\log_{10}\rho_{s}
\!\approx\!0.30\!\pm\!0.10$\,dex elevation at $z\!=\!0.85$
relative to local SPARC-equivalent halos.  In any
modified-inertia framework the surface-density scale of inferred
DM halos is set by the prevailing $\azero(z)$, so a positive
$\Delta\log_{10}\rho_{s}(z)$ at fixed baryonic content is the
DM-frame counterpart of HEAT's predicted $\azero(z)$ rise.
HEAT predicts
$\Delta\log_{10}\rho_{s}(z)\!\sim\!\log_{10}[H(z)/H_{0}]
\!=\!0.21$\,dex at $z\!=\!0.85$, within $0.9\sigma$ of the
Paper~I measurement.  Figure~\ref{fig:a0z} summarises the
$\azero(z)$ comparison and the
$\Delta\log_{10}\Sigma_{\rm DM}(z)$ counterpart.

\paragraph{Status.}  The intermediate-redshift RAR data of
\citet{Ciocan2026RAR}, anchored at $z\!=\!0$ by
\citet{Desmond2023} and \citet{McGaugh2016} and at
$z\!\sim\!0.04$ by \citet{Varasteanu2025}, jointly
\emph{exclude} the constant-$\azero$ MOND baseline at
$\Delta\chi^{2}\!\geq\!90$ on every calibration choice in
Table~\ref{tab:beta_n_degeneracy}.  The bare HEAT shape
($n\!=\!1$) is accepted by the likelihood-ratio test against
the free-$n$ alternative at $p\!=\!0.47$ on the raw cohort,
and the three-low-$z$ vs.\ four-high-$z$ predictive
cross-validation is internally consistent in both directions
at $95\%$ ($\chi^{2}_{\rm fwd}\!=\!7.89$ vs.\ $9.49$;
$\chi^{2}_{\rm back}\!=\!3.97$ vs.\ $7.81$, both under HEAT
$n\!=\!1$).  At Planck-2018 cosmology the prefactor is
$\beta\!=\!1.29\!\pm\!0.04$; this $\sim\!7\sigma$ offset from
$\beta\!=\!1$ is jointly degenerate with $H_{0}$ and with the
stellar M/L calibration, and the bare HEAT prediction is
recovered when a SH$_{0}$ES-style $H_{0}\!=\!73$ is adopted
together with a $+0.11$\,dex M/L revision (within the
per-galaxy M/L uncertainty floor of Paper~III), although the
same M/L revision over-corrects the high-$z$ data in the
predictive cross-validation.  The MOND-framework
intercept extrapolates to a $z\!=\!0$ value compatible with the
conventional Milgrom--Hubble anchor at the $\sim 1\%$ level.
Combined with the Paper~I halo-density elevation and the
Paper~II bTFR zero-evolution result discussed in
\S\ref{sec:btfr}, the trilogy provides two of three independent
observational signatures of an evolving $\azero(z)$, with the
lensed bTFR zero-evolution result remaining the principal
challenge listed in the falsification programme of
\S\ref{sec:falsify}.

\begin{figure*}[t]
  \centering
  \includegraphics[width=0.95\textwidth]{figures/fig9_a0_evolution.pdf}
  \caption{Direct $\azero(z)$ diagnostic against the MUSE-DARK
    trilogy.  \textit{Panel (a):} $\azero(z)$ values from
    \citet{Ciocan2026RAR} (Paper~III; binned values documented in
    \texttt{heat\_data/ciocan2026\_a0z.csv}) compared with the
    HEAT shape $\azero(z)\!\propto\!H(z)$.  The upper solid blue
    curve is the raw-cohort free-$K$ best fit
    ($\beta\!=\!1.29\!\pm\!0.04$, $\chi^{2}\!=\!5.9$ for $6$
    dof) -- this is the headline joint fit on Planck cosmology
    and is accepted by the likelihood-ratio test against a
    free-$n$ alternative at $p\!=\!0.47$.  The lower solid blue
    curve is the $+0.11$\,dex M/L-corrected best fit
    ($\beta\!=\!1.06\!\pm\!0.03$, $\chi^{2}\!=\!6.2$) -- this
    is shown as a calibration sensitivity, not as a preferred
    analysis (it fails the forward $3\!\to\!4$ predictive
    cross-validation at $95\%$ on the raw cohort; see text).
    The shaded blue band between them is the M/L systematic
    envelope.  The bare HEAT prediction
    $\azero(z)\!=\!\azerolocal H(z)/H_0$ with
    $\azerolocal\!=\!cH_0/(2\pi)$ ($\beta\!=\!1$) is shown as a
    dotted blue reference.  The constant-$\azero$ MOND baseline
    ($1.20\!\times\!10^{-10}\,\mathrm{m\,s^{-2}}$, dashed grey)
    is excluded at $\Delta\chi^{2}\!\geq\!90$ everywhere in the
    envelope.  The three Ciocan linear-in-$z$ fits in the DC14
    uniform, best-fit per-galaxy, and MOND frameworks are
    shown as dotted lines.  Black square: SPARC $z\!=\!0$
    RAR-fit median \citep{McGaugh2016}; black diamond:
    \citet{Desmond2023} full-Hamiltonian Monte Carlo SPARC
    re-inference (MLS interpolating function, matching
    Paper~III); red plus: \citet{Varasteanu2025} combined
    MIGHTEE-HI + high-$g$-SPARC RAR-framework value at
    $z\!\sim\!0.04$; purple circles: Ciocan
    $z\!\sim\!0.55$--$1.30$ binned RAR points (Paper~III Fig.~3,
    1-$\sigma$ error bars); purple/green/orange large markers:
    the same Ciocan global measurement at $z\!\sim\!1$ in the
    three Paper~III frameworks (Eqs.~2, D.1, E.4).  The
    upper-left inset tabulates the $\chi^{2}$ goodness-of-fit
    ladder for the six models reproduced in
    Table~\ref{tab:a0z_chi2}, on both the raw $N\!=\!7$ cohort
    and the $+0.11$\,dex M/L-corrected cohort; the annotation
    highlights the $1\%$ agreement between the Ciocan
    MOND-framework intercept ($a_0(0)\!=\!1.03\pm 0.05$) and
    the conventional Milgrom--Hubble anchor
    ($cH_0/(2\pi)\!=\!1.04$), which is the apples-to-apples
    comparison since both invert the same MOND interpolation
    function on the same kinematics.
    \textit{Panel (b):} Halo-density evolution counterpart in
    the DM frame.  Solid blue:
    $\Delta\log_{10}\,\Sigma_{\rm DM}(z)
    \!\equiv\!\log_{10}[H(z)/H_{0}]$, the HEAT-implied scaling
    of inferred DM surface density with cosmic time.
    Purple band: empirical
    $\rho_{s}\!\propto\!(1\!+\!z)^{0.54\pm 0.31}$
    fit from \citet{Ciocan2026Halo} (Paper~I).  Purple diamond:
    Paper~I MHUDF measurement at $z\!\sim\!0.85$; its offset from
    the HEAT curve is quoted in the text as a consistency-level
    comparison, not as an independent fit.\label{fig:a0z}}
\end{figure*}

\subsection{Local anchor: consistency with SPARC at
  $z\approx 0$}\label{sec:sparc}

The conventional anchor in Eq.~\ref{eq:a0} reduces to standard
MOND at $z=0$ by construction; the role of the $z=0$ comparison is
therefore to certify that the absolute-normalisation layer has not
damaged existing deep-MOND phenomenology.  Against 171 SPARC galaxies
\citep{Lelli2016} with fixed mass-to-light ratios
($\Upsilon_{\rm disk}=0.5$, $\Upsilon_{\rm bulge}=0.7\,\Msun/L_\odot$)
and no per-galaxy tuning, HEAT and constant-$\azero$ MOND are
statistically indistinguishable (paired Wilcoxon $p=0.46$);
HEAT achieves a lower $\chi^{2}_{\rm red}$ than MOND in
$98/171=57\%$ of galaxies and the same holds separately for
the Newtonian ($g_{\rm bar}>\azero$) and deep-MOND
($g_{\rm bar}<\azero$) sub-samples
(Figure~\ref{fig:sparc}).  The Milgrom--Hubble value
$\azero(0)=cH_{0}/(2\pi)=1.04\times10^{-10}\,\mathrm{m\,s^{-2}}$
is adopted from Eq.~\ref{eq:a0} rather than fitted and lies
within the $\sim 20\%$ per-galaxy scatter of the SPARC
posterior \citep{Lelli2016,Li2018}.

\begin{figure}[t]
  \centering
  \includegraphics[width=\columnwidth]{figures/sparc_chi2_red_histograms.png}
  \caption{Reduced $\chi^{2}$ distributions for 171 SPARC
    galaxies under HEAT with the conventional local anchor
    $\azero(0)=cH_0/(2\pi)$ and constant-$\azero$ MOND,
    using fixed mass-to-light ratios and no per-galaxy tuning in
    this comparison.  The distributions are
    statistically indistinguishable (paired Wilcoxon
    $p=0.46$).\label{fig:sparc}}
\end{figure}

% ======================================================================
\section{Falsification}\label{sec:falsify}
% ======================================================================

HEAT makes concrete, pre-registered predictions that distinguish
it from both $\lcdm$ and constant-$\azero$ MOND.  The thresholds
fall into two layers.  The scaling-only layer consists of the
late-type size exponent, the $\mathcal{R}(z)$ asymptote and
transition shape, and the velocity ratio; these do not depend on
the absolute value of $\azerolocal$.  The direct RAR channel has
two entries: a floated-$K$ shape test of
$\azero(z)\propto H(z)$, and an absolute-normalisation test once
$\azerolocal=cH_0/(2\pi)$ is imposed.  BTFR zero-point shifts and
the SPARC anchor complete the absolute layer.  The general
strategy of testing MOND
variants via the redshift dependence of the radial acceleration
relation was first advocated by \citet{HossenfelderMistele2018}
as the cleanest discriminator among MOND-family frameworks.

\begin{enumerate}
  \item \textbf{Velocity scaling / high-$z$ BTFR consistency.}  HEAT
    predicts $v_{\rm flat}(z)/v_{\rm flat}(0)=[\azero(z)/\azero(0)]^{1/4}$,
    with $v_{\rm flat}(z=4)/v_{\rm flat}(0)\approx1.59$.
    Constant-$\azero$ MOND predicts $1.00$ at all epochs; CEG
    \citep{Hossenfelder2017,HossenfelderMistele2018} predicts
    a substantially weaker evolution.  A sample of
    $\gtrsim 10$ BTFR points at $z=3$--$6$ with individual
    $15\%$ velocity precision discriminates all three at
    $>3\sigma$.  HEAT is falsified if the velocity normalisation
    at $z\!\sim\!4$ is measured below $1.30$.
  \item \textbf{Size-exponent test.}  HEAT predicts
    $R(z)/R_{0}\propto(1+z)^{-3/4}$ in the matter era
    (\S\ref{sec:exp}).  A measured late-type exponent
    significantly steeper than $-1$ or shallower than
    $-0.5$ between $z=0.5$ and $z=3$ would rule HEAT out.
  \item \textbf{Normalisation asymptote and transition shape.}
    Eq.~\ref{eq:Rnorm2} forces
    $\mathcal{R}_{\infty}\equiv\lim_{z\to\infty}\mathcal{R}(z)=
    \Omega_{m}^{-1/4}\!\approx\!1.34$ and the specific
    $\lcdm$-fixed transition shape between
    $\mathcal{R}(0)=1$ and $\mathcal{R}_{\infty}$.  An
    asymptote measured outside $[1.2,\,1.5]$, or a transition
    shape inconsistent with the
    $[\Omega_{m}(z)/\Omega_{m}(0)]^{1/4}$ profile at the
    percent level in Euclid/Roman data, would disfavour the
    $\azero\propto H(z)$ coupling independently of test 2.
  \item \textbf{Redshift-dependent RAR.}  At fixed
    $g_{\rm bar}$, $g_{\rm obs}\propto\azero^{1/2}$, displacing
    the RAR intercept by a factor $\sim 2.5$ at $z=4$ from
    the scaling ansatz.  A measured displacement below
    $1.3$ at $z\gtrsim 3$ rules HEAT out while remaining
    consistent with constant-$\azero$ MOND.
  \item \textbf{Direct $\azero(z\!=\!2)$ measurement.}  HEAT
    predicts $\azero(z\!=\!2)\!=\!3.03\,\azero(0)\!\approx\!3.15\!
    \times\!10^{-10}\,\mathrm{m\,s^{-2}}$.  An RAR-based
    determination of $\azero$ from a homogeneous
    $z\!\sim\!2$ rotation-curve sample (e.g.\ extending the
    \citealt{Ciocan2026RAR} pipeline to JWST/NIRSpec IFU disks)
    that places $\azero(z\!=\!2)\!<\!2.0\!\times\!10^{-10}\,
    \mathrm{m\,s^{-2}}$ at $\geq 3\sigma$ falsifies the
    absolute-anchor version of the $\azero\propto H(z)$ ansatz
    independently of the size and
    BTFR channels.
  \item \textbf{Direct $\azero(z\!=\!3)$ measurement.}  Likewise,
    HEAT predicts $\azero(z\!=\!3)\!\approx\!4.75\!\times\!
    10^{-10}\,\mathrm{m\,s^{-2}}$.  A measurement
    $<\!2.5\!\times\!10^{-10}\,\mathrm{m\,s^{-2}}$ at the same
    significance falsifies the ansatz at the high-redshift
    matter-era asymptote, complementing test~5.
  \item \textbf{Direct HI-based bTFR threshold at $z\!\sim\!1$.}
    HEAT predicts a bTFR zero-point shift
    $\Delta\log_{10}M_{b}\!=\!-\log_{10}[H(z\!=\!1)/H_{0}]\!\approx\!
    -0.255$\,dex at $z\!=\!1$.  A homogeneous direct HI-based
    bTFR sample (i.e.\ velocities measured from spatially
    resolved HI rotation in unlensed disks, rather than via
    lens-model reconstruction of strongly-lensed emission-line
    kinematics, removing the Jeanneau+2026 source-plane
    systematic) of
    $\gtrsim 30$ disks at $z\!\sim\!1$ that returns
    $|\Delta\log_{10}M_{b}|\!<\!0.05$\,dex centred on $0.00$ at
    $\geq 3\sigma$ confidence falsifies Eq.~\ref{eq:btfr-evo}.
    The current Jeanneau+2026 lensed-bTFR zero-evolution result
    ($0.00\!\pm\!0.05$\,dex) is consistent with this falsifying
    threshold and we treat it as a provisional challenge that
    awaits independent direct HI-based confirmation.
\end{enumerate}

% ======================================================================
\section{Conclusion}\label{sec:conclusions}
% ======================================================================

Adopting the scaling-only acceleration relation
$\azero(z)/\azero(0)=H(z)/H_0$ (Eq.~\ref{eq:a0scale}) yields, at
fixed baryonic mass, two zero-parameter redshift-evolution
predictions for late-type galaxy structure.  This scaling layer is
the core of the paper; the absolute Milgrom--Hubble anchor
$\azero(0)=cH_0/(2\pi)$ is used only when direct accelerations or
absolute velocities are required.  The scaling predictions combine
into the headline factorised
relation
$R_{\heat}(z)/R_{0}=(1+z)^{-3/4}\times[\Omega_{m}(z)/\Omega_{m}(0)]^{1/4}$
(Eq.~\ref{eq:Rfull}), in which exponent and normalisation enter
as separately testable factors and in which the $1/(2\pi)$
anchor cancels exactly --- so the headline test is
parameter-free under the weaker assumption
$\azero(z)\propto H(z)$ alone.  The matter-era exponent
$R\propto(1+z)^{-3/4}$ is the first galactic-scale consequence of
the acceleration scaling and is consistent at $1\sigma$
with the \citet{vanderWel2014} measurement in thousands of
galaxies.  The closed-form normalisation
$\mathcal{R}(z)=[\Omega_{m}(z)/\Omega_{m}(0)]^{1/4}$
(Eq.~\ref{eq:Rnorm2}) is a bounded, monotonic function of
redshift with a specific shape and a Planck-fixed matter-era
asymptote $\Omega_{m}^{-1/4}\!\approx\!1.34$; current data are
consistent with it at the $\sim 20$--$30\%$ level set by local
size-calibration systematics but do not yet discriminate it
from the matter-era baseline.  Upcoming Euclid/Roman data
will push the per-bin uncertainty to the percent level,
converting the present consistency check into a quantitative
test of the $\azero\propto H(z)$ coupling within the $\lcdm$
expansion history on structure scales.  An internal consistency
check is provided by the joint velocity--size cross-check on
four ALMA disks at $z\!\sim\!4.5$ (Table~\ref{tab:alma};
$N\!=\!4$, not a statistical test) and by statistical
equivalence with constant-$\azero$ MOND on SPARC at $z=0$
(\S\ref{sec:sparc}).

A covariant completion of Eq.~\ref{eq:a0scale} analogous to that of
TeVeS \citep{Bekenstein2004}, Covariant Emergent Gravity
\citep{Hossenfelder2017}, or the relativistic MOND family of
\citet{Skordis2021} is not attempted here and remains the
principal theoretical open problem.  We therefore do not
claim HEAT is confirmed; rather, we offer a zero-parameter,
scaling alternative to constant-$\azero$ MOND, with a conventional
absolute anchor for secondary tests, whose seven pre-registered
falsifiers (\S\ref{sec:falsify})
can each be resolved within the next survey cycle.

The empirical status of this article is therefore conditional
and mixed.  Current data on the headline size relation
$R(z)/R_{0}$ (\S\ref{sec:interz}) establish consistency at the
$\sim 20$--$30\%$ level set by local size-calibration systematics
but do not yet discriminate $\mathcal{R}(z)$ from the
matter-era baseline.  The recent MUSE-DARK trilogy then provides
the cleanest direct check of the acceleration-scale shape:
the Paper~III RAR-based $\azero(z)$ measurement is consistent
with the HEAT $H(z)$ form when only $K$ is floated
(\S\ref{sec:a0z}, $\chi^{2}\!=\!5.9$ for $6$ dof on the clean
$N\!=\!7$ cohort, excluding constant-$\azero$ MOND at
$\Delta\chi^{2}\!>\!190$).  The HEAT $n\!=\!1$ shape is
accepted against a free-$n$ alternative at $p\!=\!0.47$, and a
$3\!+\!4$ predictive cross-validation under HEAT $n\!=\!1$ is
internally consistent in both directions
($\chi^{2}_{\rm fwd}\!=\!7.89$ vs.\ $9.49$ and
$\chi^{2}_{\rm back}\!=\!3.97$ vs.\ $7.81$ at $95\%$).  The
joint prefactor at Planck cosmology is
$\beta\!=\!1.29\!\pm\!0.04$; this $\sim\!7\sigma$ offset from
$\beta\!=\!1$ is jointly degenerate with $H_{0}$ and stellar
M/L, and the bare HEAT prediction is reached within plausible
calibration systematics but is not the most predictively
consistent slice of that envelope.  The conventional
Milgrom--Hubble value $cH_0/(2\pi)$ remains a secondary anchor;
in the MOND-framework reduction its $z=0$ intercept matches
that anchor to within $1\%$.  The
Paper~I halo-density rise
$\rho_{s}\!\propto\!(1\!+\!z)^{0.54\pm 0.31}$ is consistent
within $0.9\sigma$ with HEAT's predicted $\Delta\log_{10}\Sigma_{\rm
DM}(z)$; the Paper~II lensed bTFR zero-evolution result at
$z\!\sim\!1$, by contrast, sits $\sim 4\sigma$ below the
HEAT-predicted shift and is the principal extant challenge to
Eq.~\ref{eq:btfr-evo},
with caveats discussed in \S\ref{sec:btfr}.  The next decisive
tests are the joint percent-level measurement of
$R(z)/R_{0}$ over $0\!\lesssim\!z\!\lesssim\!3$ from
homogeneous, mass-complete Euclid and Roman samples and the
direct HI-based bTFR test at $z\!\sim\!1$: a confirmed rise
from $\mathcal{R}(0)\!=\!1$ to
$\mathcal{R}_{\infty}\!\approx\!1.34$ with the specific
$\Omega_{\Lambda}$-sourced shape predicted by
Eq.~\ref{eq:Rnorm2} together with a confirmed
$\Delta\log_{10}M_{b}(z\!=\!1)\!\approx\!-0.26$\,dex would
provide strong support for the
$\azero\propto H(z)$ coupling on structure scales, while a
measured exponent outside $-0.75\pm 0.05$ at fixed baryonic
mass, an asymptote $\mathcal{R}_{\infty}$ inconsistent with
$\Omega_{m}^{-1/4}$ at the percent level, or a confirmed direct
HI-based bTFR zero-evolution result at $z\!\sim\!1$ would falsify
the corresponding scaling or absolute-normalisation layer specified
in \S\ref{sec:falsify}.  The full
falsification programme---including
the new $\azero(z\!=\!2)$, $\azero(z\!=\!3)$, and direct HI bTFR
thresholds---is collected in \S\ref{sec:falsify}.

% ======================================================================
\section*{Acknowledgments}
% ======================================================================

This work is independent and unfunded; it relies on publicly
available data products from the MUSE-DARK collaboration
\citep{Ciocan2026RAR,Ciocan2026Halo,Jeanneau2026}, SPARC,
MIGHTEE-HI, ALMA, JWST, and HST.

% ======================================================================
\section*{Claim ledger and sensitivity summary}
% ======================================================================

For submission clarity we separate the claims made above by evidential
status.  The primary claim is the scaling-only prediction
$R(z)/R_{0}=[H(z)/H_{0}]^{-1/2}$, equivalently
Eq.~\ref{eq:Rfull}.  It depends on the assumptions of a
mass-selected disk population, deep-MOND equilibrium scaling, and a
common local size anchor, but it does \emph{not} depend on the
absolute Milgrom--Hubble normalisation.  Current data test this claim
only at the $\sim20$--$30\%$ level because local size anchors and
mass-selection systematics dominate the error budget; a percent-level
measurement of $\mathcal{R}(z)$ with a shape inconsistent with
Eq.~\ref{eq:Rnorm2} would falsify this layer.

The second layer contains absolute-normalisation tests that require
$\azerolocal=cH_0/(2\pi)$: the flat-velocity scaling, BTFR zero-point
evolution, direct $\azero(z)$ values, and the SPARC local comparison.
These are not used to establish the headline size relation.  The
MUSE-DARK RAR comparison is therefore reported in two forms: the
free-$K$ shape test of $\azero(z)\propto H(z)$, and the conventional
$cH_0/(2\pi)$ anchor as a secondary reference.  The shape test is
backed by six independent robustness exercises whose outputs are
collected in \texttt{publication/a0\_robustness\_tests.py}:
(A) a joint $(\beta, n)$ MCMC posterior; (B) a $9\!\times\!5$
cosmology grid scan in $(H_{0}, \Omega_{m})$; (C) the
likelihood-ratio test for $n\!=\!1$ vs.\ free $n$; (D) a
$10^{4}$-draw bootstrap of the $(\beta, n)$ fit; (E) a
bidirectional $3\!+\!4$ predictive cross-validation under HEAT
$n\!=\!1$; and (F) a smooth $z$-dependent M/L profile
sensitivity.  The Jeanneau+2026 lensed bTFR zero-evolution result
is the strongest current challenge to the absolute-normalisation
layer and is explicitly listed among the falsification thresholds.

The ALMA $z\!\sim\!4.5$ velocity--size sample and the JWST compact
systems are supporting consistency checks only.  Their small sample
size, lensing/merger caveats, gas-fraction uncertainties, and
non-equilibrium dynamics prevent them from being treated as a
standalone statistical confirmation.  They are retained because they
probe the same redshift scaling with different powers of $\azero$ and
therefore define useful near-term tests for larger samples.

No covariant completion, gravitational-lensing calculation, CMB
perturbation theory, or cluster solution is claimed in this article.
Those observables require additional dynamics beyond the scaling ansatz
of Eq.~\ref{eq:a0scale}; failure in any of those sectors would constrain
candidate completions rather than being hidden by a parameter adjustment
inside the present letter model.

% ======================================================================
\section*{Data and code availability}

All code, data catalogues, and the \LaTeX\ source of this article
are publicly available at
\url{https://github.com/bogdanenciu/HEAT-paper} and archived with
a permanent Digital Object Identifier at
\dataset[10.5281/zenodo.19709780]{https://doi.org/10.5281/zenodo.19709780}.
The compiled preprint PDF is separately archived at
\dataset[10.5281/zenodo.19710286]{https://doi.org/10.5281/zenodo.19710286}
(same title; concept DOI tracks the latest PDF version).
The MUSE-DARK trilogy public catalogues
(\citealt{Ciocan2026Halo}~Paper~I MHUDF photometry,
\citealt{Ciocan2026RAR}~Paper~III binned $\azero(z)$, and the
\citealt{Varasteanu2025} MIGHTEE-HI anchor) used as overlays in
Figs.~\ref{fig:size}, \ref{fig:btfr}, and~\ref{fig:a0z} are
bundled in \texttt{heat\_data/} with column-level provenance.  The
figure set of this article can be regenerated from a clean
checkout with
\texttt{python run\_all\_heat\_tests.py all --agg}.

\software{NumPy \citep{Harris2020},
          SciPy \citep{Virtanen2020},
          Matplotlib \citep{Hunter2007}}

% ======================================================================
% Bibliography
% ======================================================================
\bibliographystyle{aasjournal}

\begin{thebibliography}{99}
\bibitem[Allen et al.(2017)]{Allen2017}
  Allen, R.~J., Kacprzak, G.~G., Glazebrook, K., et al.\ 2017, ApJL,
  834, L11
\bibitem[Allen et al.(2025)]{Allen2025}
  Allen, N., Oesch, P.~A., Toft, S., et al.\ 2025, A\&A,
  698, A30
\bibitem[Bak \& Rey(2000)]{BakRey2000}
  Bak, D., \& Rey, S.-J. 2000, Classical and Quantum Gravity,
  17, L83
\bibitem[Bekenstein(2004)]{Bekenstein2004}
  Bekenstein, J.~D. 2004, PRD, 70, 083509
\bibitem[Bunker et al.(2023)]{Bunker2023}
  Bunker, A.~J., et al.\ 2023, A\&A, 677, A88
\bibitem[Cai \& Kim(2005)]{CaiKim2005}
  Cai, R.-G., \& Kim, S.~P. 2005, JHEP, 02, 050
\bibitem[Carniani et al.(2024)]{Carniani2024}
  Carniani, S., et al.\ 2024, Nature, 633, 318
\bibitem[Cathey et al.(2024)]{Cathey2024}
  Cathey, J., Gonzalez, A.~H., Lower, S., et al.\ 2024, ApJ,
  967, 11
\bibitem[Ciocan et al.(2026a)]{Ciocan2026Halo}
  Ciocan, B.~I., Bouch\'e, N.~F., Fensch, J., et al.\ 2026a,
  A\&A, in press, arXiv:2506.19721 (MUSE-DARK Paper~I:
  intermediate-$z$ dark-matter halo properties)
\bibitem[Ciocan et al.(2026b)]{Ciocan2026RAR}
  Ciocan, B.~I., Bouch\'e, N.~F., Fensch, J., et al.\ 2026b,
  A\&A, in press, arXiv:2604.22613 (MUSE-DARK Paper~III:
  evolution of the radial acceleration relation at
  intermediate redshifts)
\bibitem[Conselice(2014)]{Conselice2014}
  Conselice, C.~J. 2014, ARA\&A, 52, 291
\bibitem[Desmond(2023)]{Desmond2023}
  Desmond, H. 2023, MNRAS, 526, 3342 (arXiv:2303.11314; full
  Hamiltonian Monte Carlo joint inference of the underlying
  SPARC RAR under the MLS interpolating function, yielding
  $\azero(0)\!=\!1.19\!\pm\!0.04\,(\mathrm{stat})
  \!\pm\!0.09\,(\mathrm{sys})\!\times\!10^{-10}\,
  \mathrm{m\,s^{-2}}$)
\bibitem[Duncan et al.(2019)]{Duncan2019}
  Duncan, K., Conselice, C.~J., Mortlock, A., et al.\ 2019, ApJ,
  876, 110
\bibitem[Famaey \& McGaugh(2012)]{FamaeyMcGaugh2012}
  Famaey, B., \& McGaugh, S.~S. 2012, Living Reviews in Relativity,
  15, 10
\bibitem[Gibbons \& Hawking(1977)]{GibbonsHawking1977}
  Gibbons, G.~W., \& Hawking, S.~W. 1977, PRD, 15, 2738
\bibitem[Harris et al.(2020)]{Harris2020}
  Harris, C.~R., et al.\ 2020, Nature, 585, 357
\bibitem[Hawking(1975)]{Hawking1975}
  Hawking, S.~W. 1975, Communications in Mathematical Physics,
  43, 199
\bibitem[Hossenfelder(2017)]{Hossenfelder2017}
  Hossenfelder, S. 2017, PRD, 95, 124018
\bibitem[Hossenfelder \& Mistele(2018)]{HossenfelderMistele2018}
  Hossenfelder, S., \& Mistele, T. 2018, International Journal of
  Modern Physics D, 27, 1847010
\bibitem[Hunter(2007)]{Hunter2007}
  Hunter, J.~D. 2007, Computing in Science and Engineering, 9, 90
\bibitem[Jeanneau et al.(2026)]{Jeanneau2026}
  Jeanneau, R., Bouch\'e, N.~F., Ciocan, B.~I., et al.\ 2026,
  A\&A, in press, arXiv:2603.28856 (MUSE-DARK Paper~II:
  lensed baryonic Tully--Fisher relation at $z\!\sim\!1$)
\bibitem[Labb\'e et al.(2023)]{Labbe2023}
  Labb\'e, I., et al.\ 2023, Nature, 616, 266
\bibitem[Lelli et al.(2016)]{Lelli2016}
  Lelli, F., McGaugh, S.~S., \& Schombert, J.~M. 2016, AJ, 152, 157
\bibitem[Lelli et al.(2019)]{Lelli2019}
  Lelli, F., McGaugh, S.~S., Schombert, J.~M., Desmond, H., \&
  Katz, H. 2019, MNRAS, 484, 3267
\bibitem[Li et al.(2018)]{Li2018}
  Li, P., Lelli, F., McGaugh, S.~S., \& Schombert, J.~M. 2018, A\&A,
  615, A3
\bibitem[McGaugh et al.(2016)]{McGaugh2016}
  McGaugh, S.~S., Lelli, F., \& Schombert, J.~M. 2016, PRL, 117, 201101
\bibitem[McGaugh et al.(2018)]{McGaugh2018}
  McGaugh, S.~S., Li, P., Lelli, F., \& Schombert, J.~M. 2018,
  Nature Astronomy, 2, 924
\bibitem[Milgrom(1983)]{Milgrom1983}
  Milgrom, M. 1983, ApJ, 270, 365
\bibitem[Milgrom(1999)]{Milgrom1999}
  Milgrom, M. 1999, Physics Letters A, 253, 273
\bibitem[Mo, Mao, \& White(1998)]{MoMaoWhite1998}
  Mo, H.~J., Mao, S., \& White, S.~D.~M. 1998, MNRAS, 295, 319
\bibitem[Rizzo et al.(2020)]{Rizzo2020}
  Rizzo, F., Vegetti, S., Powell, D., et al.\ 2020, Nature, 584, 201
\bibitem[Rodrigues et al.(2018)]{Rodrigues2018}
  Rodrigues, D.~C., Marra, V., del Popolo, A., \& Davari, Z. 2018,
  Nature Astronomy, 2, 668
\bibitem[Roman-Oliveira et al.(2023)]{RomanOliveira2023}
  Roman-Oliveira, F., et al.\ 2023, MNRAS, 521, 1045
\bibitem[Roman-Oliveira et al.(2024)]{RomanOliveira2024}
  Roman-Oliveira, F., et al.\ 2024, A\&A, 687, A35
\bibitem[Shen et al.(2003)]{Shen2003}
  Shen, S., Mo, H.~J., White, S.~D.~M., et al.\ 2003, MNRAS, 343, 978
\bibitem[Skordis \& Z{\l}o\'snik(2021)]{Skordis2021}
  Skordis, C., \& Z{\l}o\'snik, T. 2021, PRL, 127, 161302
\bibitem[Somerville et al.(2018)]{Somerville2018}
  Somerville, R.~S., et al.\ 2018, MNRAS, 473, 2714
\bibitem[Tacconi et al.(2018)]{Tacconi2018}
  Tacconi, L.~J., Genzel, R., Saintonge, A., et al.\
  2018, ApJ, 853, 179
\bibitem[{\"U}bler et al.(2017)]{Ubler2017}
  {\"U}bler, H., F\"orster Schreiber, N.~M., Genzel, R., et al.\
  2017, ApJ, 842, 121
\bibitem[Unruh(1976)]{Unruh1976}
  Unruh, W.~G. 1976, PRD, 14, 870
\bibitem[van der Wel et al.(2014)]{vanderWel2014}
  van der Wel, A., et al.\ 2014, ApJ, 788, 28
\bibitem[V\u{a}ra\c{s}teanu et al.(2025)]{Varasteanu2025}
  V\u{a}ra\c{s}teanu, T., Ponomareva, A., Maddox, N., et al.\
  2025, MNRAS, in press (MIGHTEE-HI direct $\azero$
  measurement at $z\!<\!0.08$)
\bibitem[Virtanen et al.(2020)]{Virtanen2020}
  Virtanen, P., et al.\ 2020, Nature Methods, 17, 261
\end{thebibliography}

\end{document}
